%
% ---------- header -----------------------------------------------------------
%
% project       kaneton
%
% license       kaneton
%
% file          /home/mycure/kaneton/view/book/development/development.tex
%
% created       julien quintard   [mon may 14 19:56:45 2007]
% updated       julien quintard   [fri may 25 19:52:31 2007]
%

%
% path
%

\def\path{../..}

%
% template
%

%
% ---------- header -----------------------------------------------------------
%
% project       kaneton
%
% license       kaneton
%
% file          /home/mycure/kaneton/view/template/book.tex
%
% created       julien quintard   [wed may 16 18:16:50 2007]
% updated       julien quintard   [sun jun  3 12:00:37 2007]
%

%
% class
%

\documentclass[10pt,a4wide]{book}

%
% packages
%

\usepackage[english]{babel}
\usepackage[T1]{fontenc}
\usepackage{a4wide}
\usepackage{fancyheadings}
\usepackage{multicol}
\usepackage{indentfirst}
\usepackage{graphicx}
\usepackage{color}
\usepackage{xcolor}
\usepackage{verbatim}

\usepackage{aeguill}

\usepackage[Lenny]{\path/package/fncychap}

\pagestyle{fancy}

%
% -
%

\renewcommand{\-}{\vspace{\parskip}}

%
% toc
%

\newcommand{\toc}
  {
    \tableofcontents
    \setlength{\footrulewidth}{0.3pt}
    \setlength{\parindent}{0.3cm}
    \setlength{\parskip}{2ex plus 0.5ex minus 0.2ex}
  }

%
% logos
%

\newcommand{\logos}
  {
    \begin{center}
      \includegraphics[scale=0.8]{\path/logo/kaneton.pdf}
    \end{center}
  }

%
% colors
%

\definecolor{functioncolor}{rgb}{0.40,0.00,0.00}
\definecolor{commandcolor}{rgb}{0.00,0.00,0.40}
\definecolor{verbatimcolor}{rgb}{0.00,0.40,0.00}
\definecolor{noticecolor}{rgb}{0.87,0.84,0.02}

%
% function
%

\newcommand\function[3]{
  \begin{tabular}{p{0.2cm}p{13.8cm}}
  & {\color{functioncolor}\textbf{#1}}#2
  \end{tabular}

  \begin{tabular}{p{1cm}p{13cm}}
  & #3
  \end{tabular}}

%
% align
%

\newcommand\align[1]{
  \\ & \hspace{#1}}

%
% argument
%

\newcommand\argument[1]{\textit{#1}}

%
% command
%

\newcommand\command[2]{
  \begin{tabular}{p{0.2cm}p{13.8cm}}
  & {\color{commandcolor}\textbf{#1}}
  \end{tabular}

  \begin{tabular}{p{1cm}p{13cm}}
  & #2
  \end{tabular}}

%
% notice
%

\newcommand\notice[1]{
  {\color{noticecolor}\textbf{Notice}}

  \begin{tabular}{p{0.2cm}p{13.8cm}}
  & #1
  \end{tabular}}

%
% subsubsubsection
%

\newcommand\subsubsubsection[1]{\textbf{#1}}

%
% example
%

\newcommand\example[1]{
  \textit{Example:}

  \begin{tabular}{p{0.2cm}p{13.8cm}}
  & \textit{#1}
  \end{tabular}}

%
% warning XXX
%

\renewcommand{\familydefault}{\sfdefault}

%
% verbatim stuff
%

\makeatletter

\renewcommand{\verbatim@font}
  {\ttfamily\footnotesize\selectfont}

\def\verbatim@processline{
  \hskip15ex{\color{verbatimcolor}\the\verbatim@line}\par
}

\makeatother

%
% header
%

\rhead{}
\rfoot{\scriptsize{The kaneton microkernel project}}

\date{\scriptsize{\today}}


%
% header
%

\lhead{\scriptsize{The kaneton microkernel :: development book}}
\rhead{}

%
% title
%

\title{The kaneton microkernel :: development book
       \logos}

%
% authors
%

\author{\small{Julien Quintard}}

%
% document
%

\begin{document}

%
% title
%

\maketitle

%
% --------- text --------------------------------------------------------------
%

This document describes how people can contribute to the kaneton microkernel
project and the rules they have to follow.

This document must be read by everyone involved in the kaneton microkernel
project.

All the kaneton documents are available on
the official website
  \footnote{http://www.kaneton.org}.

%
% toc
%

\tableofcontents

%
% chapters
%

%
% ---------- header -----------------------------------------------------------
%
% project       kaneton
%
% license       kaneton
%
% file          /home/mycure/kaneton/view/book/development/introduction.tex
%
% created       julien quintard   [thu may 17 12:46:30 2007]
% updated       julien quintard   [fri jun 15 11:48:28 2007]
%

%
% ---------- introduction -----------------------------------------------------
%

\chapter{Introduction}
\label{chapter:environment}

In this chapter, the kaneton microkernel project is briefly introduced
in order to emphasize some of its characteristics that makes contribution
specific in several ways.

\newpage

%
% ---------- text -------------------------------------------------------------
%

kaneton is a educational purpose microkernel project. This project aims
at providing a very clear, commented and maintainable microkernel source
code in order to allow people interested in operating systems internals
to look at the source code and understand it very quickly.

The kaneton project is basically composed of the source code of the
microkernel itself, scripts to perform complex tasks and various documents
from design papers to lecture materials.

The most important thing to remember is that the whole project is intended
to be understood as well as possibly maintained by everyone. As a result,
contributions must comply with the level of clarity expected by the project.

These rules are discussed in this paper in order to inform every new
contributor of what makes a good contribution.

The remaining of this document is organised as follows. \textit{Chapter
\ref{chapter:history}} draws the history of the kaneton microkernel project.
\textit{Chapter \ref{chapter:source tree}} introduced the kaneton project
organisation through the source code hierarchy. Next, \textit{Chapter
\ref{chapter:community}} describes how a contributor should behave in a
development community. \textit{Chapter \ref{chapter:rules}} introduces the
general rules which apply to any context around the kaneton project. Then,
the tools inherent to the kaneton project are listed in \textit{Chapter
\ref{chapter:tools}} with some guidelines about how to use them properly.
\textit{Chapter \ref{chapter:languages}} explicitly describes languages rules
informing the developer of the coding style to respect. \textit{Chapter
\ref{chapter:people}} draws a list of the people in charge for the different
parts and tools of the project. Finally, \textit{Chapter
\ref{chapter:licenses}} contains information about the licenses related to
the kaneton microkernel project.

%
% ---------- header -----------------------------------------------------------
%
% project       kaneton
%
% license       kaneton
%
% file          /home/mycure/kaneton/view/book/development/source-tree.tex
%
% created       julien quintard   [thu may 17 22:41:36 2007]
% updated       julien quintard   [sun may 20 18:09:37 2007]
%

%
% ---------- source tree ------------------------------------------------------
%

\chapter{Source Tree}

In this chapter we will briefly describe the kaneton microkernel project
source tree.

\newpage

%
% ---------- text -------------------------------------------------------------
%

The kaneton microkernel reference source tree looks like the following
listing:

\begin{verbatim}
cheat/
environment/
export/
configure/
kaneton/
library/
license/
tool/
test/
transcript/
view/
\end{verbatim}

% cheat/

\subsection*{cheat/}

Since the kaneton microkernel is implemented by students, the kaneton
people need to check whether students are cheating by re-using parts of
previous years projects or other kernel source code available on the
\textit{Internet}.

To avoid cheating, kaneton people developed a software checking for
commonalities between different source codes.

The \textit{cheat/} kaneton subdirectory contains everything necessary
to perform such tasks including the software, the different source codes
found on the internet, the kaneton students' implementation from the
previous years etc.

% environment/

\subsection*{environment/}

This directory contains everything necessary to the kaneton development
environment.

The kaneton development environment allows different developers to
interact on the development of the same microkernel in a pretty easy way.

The development environment aims at providing developers to possibility to
work in a collaborative manner without interfering with each other. These
developers are likely to run different operating systems on different
microprocessors. In addition, the kaneton microkernel can be targeted for
different microprocessor architectures. The development environment was
introduced to cope with these combinaisons by providing profiles, each
profile describing the behaviour of a component: underlying operating system,
target architecture, user-specific stuff etc.

As a result, each developer can use a different operating system and
microprocessor architecture with his own specific compiling flags, kaneton
parameters etc. without modifying another developer's configuration.

% export/

\subsection*{export/}

The \textit{export/} directory contains scripts used to generate a kaneton
tarball in order to be distributed to the students at the beginning of the
kaneton project.

Indeed, these scripts rearrange the kaneton hierarchy hidding some
important directories the students do not need to be aware of.

Moreover some source code parts are removed since the students then have to
rewrite these piece of code.

% configure/

\subsection*{configure/}

This directory contains everything necessary for configuring its own
kaneton microkernel development environment through the compiling process
to the boot system.

Any new contributor should first look at this directory.

% kaneton/

\subsection*{kaneton/}

This directory is the most important of the project since it contains
the whole microkernel source code.

The directory is composed of three important subdirectories: \textit{core/},
\textit{platform} and \textit{architecture}. These subdirectories are described
next.

% kaneton/core/

\subsection*{kaneton/core/}

This directory contains the kaneton core source code.

The directory is divided as shown below:

\begin{verbatim}
as/
region/
sched/
segment/
set/
task/
thread/
[...]
\end{verbatim}

Each directory represents a kaneton core manager. For more information on
the kaneton core, please refer to the appropriate document:
\textit{The kaneton microkernel :: reference implementation}

% kaneton/platform/

\subsection*{kaneton/platform/}

This directory contains everything in relation with what the kaneton
microkernel project calls a \textbf{platform}. The platform represents the
board supporting the devices: microprocessor, memory, peripherals etc.

This directory obviously contains subdirectories for each platform
supported by the kaneton microkernel.

% kaneton/architecture

\subsection*{kaneton/architecture}

The \textit{architecture} directory contains the source-code related to
the microprocessor architecture.

This directory is composed of subdirectories, each one representing an
architecture: \textit{ia32}, \textit{mips64} etc.

% library/

\subsection*{library/}

This directory contains the libraries used either by the kaneton microkernel
itself or by the kaneton microkernel servers. This directory especially
contrains the standard \textit{kaneton C library}.

% license/

\subsection*{license/}

This directory contains the licenses used for any program or document
in relation with the kaneton microkernel project.

Each student has to read and accept the kaneton license before implementing
or even using the kaneton microkernel project.

% tool/

\subsection*{tool/}

This directory contains additional files used by the kaneton development
environment.

% test/

\subsection*{test/}

Since the kaneton microkernel is used as a material for operating system
courses, the kaneton microkernel reference which is the basis of students'
work is to be extremely reliable.

The kaneton project therefore contains a set of tools in order to validate
the kaneton reference implementation behaviour. These tools are also used
for evaluating the correctness of the students' implementation.

The \textit{test/} directory so contains the set kaneton scripts and tests
for validating a kaneton microkernel implementation.

% transcript/

\subsection*{transcript/}

This directory contains real-time recorded sessions. These sessions can be
replayed showing how to use the kaneton microkernel project and more
specifically the kaneton development environement.

% view/

\subsection*{view/}

This directory contains all the kaneton documents including kaneton
administrative documents, exams documents, lectures materials, kaneton papers
and books etc.

Additionally, scripts are provided in order to very easily build and
display these documents.

%
% ---------- header -----------------------------------------------------------
%
% project       kaneton
%
% license       kaneton
%
% file          /home/mycure/kaneton/view/book/development/community.tex
%
% created       julien quintard   [sun may 20 18:08:17 2007]
% updated       julien quintard   [tue may 20 17:03:59 2008]
%

%
% ---------- community --------------------------------------------------------
%

\chapter{Community}
\label{chapter:community}

This chapter discusses what is a community and how contributors must integrate
the kaneton community.

\newpage

%
% ---------- text -------------------------------------------------------------
%

kaneton can obviously be considered as an open source community although the
produced soure code is actually not open source.

Driving an open source community is complicated since people have different
personal goals at working on a free project. Some people contribute for
the knowledge, other for building the next generation system, other to
provide free open source softwares, other to become famous \etc{}

kaneton is a community driven microkernel that acts with the best interest of
the students at heart. Rules and regulations that keep the project moving
forward are fundamental even if the size of the kaneton community is
relatively small, for now.

Indeed, the main objective of the kaneton project remains to be as
understandable as possible in order to lead students to implement parts of
it very quickly.

The remaining of this chapter draws a list of rules contributors must agree
to respect.

% objective

\subsubsection{Objective}

kaneton aims at providing a powerful, understandable and maintainable
microkernel. This objective must be kept in mind of every contributor
since many design and implementation were/are/will be made according to
this precise objective.

Note that the kaneton microkernel does not intend to be a desktop operating
system nor an as optimised as the Linux operating system. Every contributor
should be well-aware of that in order to avoid behaviours stating that a
feature is fundamental or useless for performance concerns, for instance.

This rule does not prohibit people to suggest ideas but instead regulates
behaviours of people who wants to change major design and/or implementation
choices for bad reasons.

% behaviour

\subsubsection{Behaviour}

Open source projects does not mean constraint-free projects. The kaneton
people, whilst being relatively young, try to act for the project's good
by behaving remarkably in the kaneton community.

Therefore, contributors are asked to do the same by avoiding some bad/young
behaviours.

\begin{enumerate}
  \item
    Follow the rules. People who do not respect these rules could be banned
    from the kaneton project.
  \item
    Avoid the \textit{cowboy} behaviour consisting for a contributor to
    implement a feature without discussing about its usefulness with the
    community first. Another effect of this behaviour can be to distract
    the contributors from its major focuses.
  \item
    Always act and think in the project interest rather than your personal
    interest.
  \item
    Respect the other kaneton people, especially the ones who have worked
    on this project for a long time and who made this whole project possible.
    When people disagree, they are asked to do it respectfully.
  \item
    Take your responsability when you realise that you did something wrong:
    insults, mistakes in an implemented feature \etc{}
  \item
    \textit{``The Perfect is the Enemy of the Good!''}: even nice
    contributors can unintentionally do bad things by being perfectionists
    and/or too much into the project and/or obsessed with process.
  \item
    ... \textit{Politness}, \textit{Respect}, \textit{Trust} and
    \textit{Humility} are the key qualities that make a good contributor in
    any community.
\end{enumerate}

% communication

\subsubsection{Communication}

The communication mainly takes two forms in the kaneton microkernel project:
the \name{mailing-list} for internal communication and the \name{kaneton
website} for external communication. The \name{Developers Wiki} is another
source of communication as well as the commit logs \etc{}

The rules related to these tools are described in \reference{Chapter
\ref{chapter:tools}} and will therefore not be discussed here.

Every contributor must take the time to communicate in the mailing-list as well
as through the public website. Indeed, kaneton people must, frequently,
briefly describe what they are working on in order to inform the other
contributors who are not aware of everyone's current work. Note that these
kind of messages are very different from messages generated by repository
commits. Indeed, while these commit messages indicate a modification, they
do not describe the whole work behind them.

Additionally, contributors can communicate informing the kaneton community of
their unavailability for the next two months, for instance. This behaviour
allows people to be aware that some tasks will not be done because the
contributors in charge of it cannot work at this moment.

Althoug people are highly welcomed to communicate, some rules apply to
avoid further problems.

First, any new contributor should obviously read the kaneton documents and
then go through the mailing-list archives before asking anything which has
already been discussed, unless the developer knows exactly what he/she is
talking about. Indeed, asking too many questions about the source code is a
form of disrespect to the other contribors. Moreover, many things can be
found out just looking at the kaneton documents and/or source code.

Although people are asked to communicate, people are also asked to act
respectully. Contributors should not respond to every message in every
discussion, this is a ridiculous behaviour. Instead, every developer should
carefully read the discussion, think about its response and then write a clear
message stating his point of view, ideas \etc{}

Depending on the contributor status, reading the mailing-list frequently is
absolutely fundamental as some people rely on other contributors' decisions,
advices \etc{}

Finally, the mailing-list must be considered as the official internal
communication medium. If contributors previously had a private conversation,
in real-life or on \name{IRC} for instance, the discussion must be reported
on the mailing-list so that everyone can take these new ideas into account.

% work

\subsubsection{Work}

Working on the kaneton microkernel project does not imply low-level programming
all the time. Indeed, the kaneton project is composed of two parts: the
kaneton microkernel research project and the kaneton educational project.

Although the microkernel research project requires highly skilled programmers,
it also needs documentations and some tools for performing important tasks as
diverse as generating the prototypes, testing the microkernel behaviour,
generating the documentation \etc{}

The educational project essentially needs documentation, lecture materials
and tools for managing the project: testing the students' implementation,
checking if some students cheated and many others.

This means that kaneton people must contribute to every type of task
that need to be done. Also, contributors are asked to well document
any work they have done including source code comments but also through
kaneton official documents which are then made available on the website.

% supervisor

\subsubsection{Supervisor}

A supervisor is attached to each new contributor for a certain, not fixed,
period of time. The role of this supervisor is to advise, correct and
encourage the newcomer so that it integrates well the kaneton community.

The supervisor will be someone having a well understanding of the new
contributor's project. A high level of communication is expected between
the supervisor and the contributor. A direct phone communication is highly
recommanded through \name{Skype} for instance.

% trust

\subsubsection{Trust}

A new contributor joining the kaneton project must acquire the trust of
the community. Therefore, the contributor first does not get any access
to restricted tools and must submit patches to its supervisor who review
them, taking care of advising the newcomer of its mistakes.

At the end of the test period, the community decides whether the contributor
is helpful to the project or not. Then, the contributor is either granted
of full access to kaneton tools or eliminated from the project.

%
% ---------- header -----------------------------------------------------------
%
% project       kaneton
%
% license       kaneton
%
% file          /home/mycure/kaneton/view/book/development/tools.tex
%
% created       julien quintard   [sun may 20 14:48:11 2007]
% updated       julien quintard   [sun may 20 15:22:11 2007]
%

%
% ---------- tools ------------------------------------------------------------
%

\chapter{Tools}

This chapter describes every tool kaneton contributors use on a daily-basis.

\newpage

%
% ---------- text -------------------------------------------------------------
%

%
% internal
%

\section{Internal}

%
% ---------- header -----------------------------------------------------------
%
% project       kaneton
%
% license       kaneton
%
% file          /home/mycure/kaneton/view/book/development/environment.tex
%
% created       julien quintard   [sun may 20 14:49:26 2007]
% updated       julien quintard   [thu jun 12 11:06:22 2008]
%

%
% ---------- environment ------------------------------------------------------
%

\subsection{Environment}
\label{section:environment}

Over the years, the kaneton microkernel evolved, starting with a very simple
introduction to low-level programming and finally to a complete microkernel
development.

kaneton people wanted to lead students to a complete microkernel development
to finally introduce distributed computing. This would not have been possible
if students had to build an entire development environment because developing
such an environment is a whole project by itself.

As a result, kaneton people decided to provide students a complete development
environment. The kaneton development environment is composed of make files,
python scripts and configuration files. This development environment can be
considered as one of the major kaneton tools since contributors use it
everytime.

The kaneton development environment aims at providing an easy and portable
way for managing the kaneton microkernel project from a development point
of view. Therefore, the kaneton environment provides everything necessary
for compiling, assembling, \etc{} These tasks highly rely on the underlying
running operating system as well as on the kaneton microkernel's target
microprocessor. Moreover, the user could need to redefine some behaviours
depending on its personal operating system configuration to use a specific
C compiler for instance.

The kaneton development environment provides a layered organisation of
profiles, each profile defining variables and functions used by the final
environment engine. The goal of the layered model is to allow layers to
override the definitions of lower-level layers.

%
% profiles
%

\subsubsection{Profiles}

The development environment is composed of profiles including a \name{host}
profile which describes the behaviour of the underlying operating system, a
\name{kaneton} profile which parameters the kaneton microkernel and a
\name{user} profile which permits the user to redefine lower-level layers'
declarations.

These profiles eventually hold sub-profiles which define variables and
functions. These actual profiles are accessed according to user-defined
shell variables.

% host

\subsubsubsection{Host}

The \name{host} profile essentially describes how to perform basic tasks:
compile, assemble, change the current directory, display a message \etc{} These
tasks rely on the operating system currently running as well as on the target
processor which kaneton will be built for. For these reasons, there are
several host sub-profiles.

Let us suppose a developer is running a \name{Linux} operating system and
that kaneton will be built for running on a \name{PowerPC} microprocessor. In
such a case, the C compiler program will be different depending on the
microprocessor \name{Linux} is running on. Indeed, if Linux is running on
a \name{PowerPC} microprocessor, then using the default compiler should
produce \name{PowerPC} object files. This is well-known to be the common
compiling way. On the other hand, if \name{Linux} is running on a
different microprocessor, then a cross-compiler must be used to produce
binary objects targeting a specific different microprocessor architecture,
the \name{PowerPC} architecture in our example.

To avoid this issue, a \name{host} sub-profile name is composed of two parts
separated by a slash. The first part is the name of the operating system and
the latter is a pair source/target processors separated by a period. For
example, \name{linux/ia32.ppc} names a \name{host} profile representing a
\name{Linux} operating system on a \name{Intel 32-bit} microprocessor
which aims at building a kaneton microkernel for a \name{PowerPC}
target architecture. Needless to say that \name{linux/ia32.ia32} represents
a non cross-compiling environment.

To avoid configuration duplications, it is common to see the configuration
file of a host sub-profile to include files of the parent directory as
shown below:

\begin{verbatim}
  linux/
    linux.desc
    linux.conf
    linux.mk
    linux.py
    ia32.ia32/
      educational -> .
      optimised -> .
      smp -> .
      ia32.desc
      ia32.conf
      ia32.mk
      ia32.py
    ia32.mips64/
      mips64.desc
      mips64.conf
      mips64.mk
      mips64.py
\end{verbatim}

Note that the files \location{linux.*} are not directly included by the
development environment engine since \name{linux} is not a valid host
profile name.

Two host profiles are illustrated here. The first one is named
\name{linux/ia32.ia32} while the second's name is \name{linux/ia32.mips64}.

For example, the \name{linux/ia32.mips64} \name{host} profile represents a
\name{Linux} operating system running on a \name{Intel 32-bit}
microprocessor while kaneton is built for a \name{MIPS 64-bit} target
architecture. This profile is likely to include the \location{linux.*} files
of the parent directory since there are not much difference between all the
\name{linux/*.*} \name{host} profiles. However, such a profile will
certainly redefine the binary paths of the C compiler, linker \etc{} in order
to produce \name{MIPS 64-bit} binary objects.

To conclude, the \name{host} sub-profile is accessed by the following
construct:

\begin{verbatim}
  profile/host/${KANETON_HOST}/${KANETON_ARCHITECTURE}
\end{verbatim}

With, for instance, the following values:

\begin{verbatim}
  KANETON_HOST = linux/ia32
  KANETON_ARCHTECTURE = ia32/educational
\end{verbatim}

Note that the possibility to include files in the configuration syntax allows
very similar profiles to share a huge amount of definitions.

% kaneton

\subsubsubsection{kaneton}

The \name{kaneton} profile is composed of four sub-profiles: \name{core},
\name{platform}, \name{architecture} and \name{libc}.

The \name{core} sub-profile contains variables for parameterizing the
kaneton core. The \name{platform} and \name{architecture} sub-profiles
focus on the configuration of the platform- and architecture-dependent code
of the kaneton microkernel, respectively. The \name{libc} sub-profile
contains parameters for the \name{C Library} used in the kernel.

The user-defined shell variables \code{\$\{KANETON\_PLATFORM\}} and
\code{\$\{KANETON\_ARCHITECTURE\}} are used to address the \name{platform}
and \name{architecture} sub-profiles, respectively.

% user

\subsubsubsection{User}

Let us suppose that a developer would like the kaneton microkernel to
use a specific memory management entirely based on a \name{Slab Allocator}
and with all microprocessor optimisations enabled. These user-specific
configurations are actually allowed by the \name{user} profile.

The user-defined shell variable \code{\$\{KANETON\_USER\}} defines the name
of the \name{user} profile. This profile contains user-specific
configurations allowing a contributor to overwrite lower-level layer defintions
in order to specialise a behaviour.

The kaneton project also provides a tool allowing users to configure
their development environment. This tool is named \name{configure} and is
available from the kaneton project root directory. For more information about
this tool, please refer to \reference{Section \ref{section:configure}}.


%
% requirements
%

\subsubsection{Requirements}

The whole kaneton development environment needs exactly two fundamental tools
to work. The first one is \name{GNU make}, used to build powerful make files,
and the second one is \name{Python}, used to write portable scripts. If an
operating system has these two tools, then kaneton can certainly be developed
on it.

As said previously, the user has to specify some shell variables which are
used by the kaneton development environment engine. These variables are
described below:

\begin{itemize}
  \item
    \code{\$\{KANETON\_USER\}}: the name of the kaneton developer.

    A \name{user} profile name must be composed of the first name, a period
    and finally, the last name of the developer.
  \item
    \code{\$\{KANETON\_HOST\}}: the name of the host which is composed of
    a couple operating system/microprocessor.
  \item
    \code{\$\{KANETON\_PYTHON\}}: the path of the python binary.

    This path is required since the very first scripts which set up the
    configured environment are written in \name{Python}.
  \item
    \code{\$\{KANETON\_PLATFORM\}}: the name of the target platform.
  \item
    \code{\$\{KANETON\_ARCHITECTURE\}}: the name of the target architecture.
\end{itemize}

Note that once the configured environment is set up, these variables are
no longer used by the kaneton environment engine. Indeed, instead, the
kaneton environment operations are based on the \name{host} profile on
which rely the configured environment.

The profiles names must all be in lowercase letters. Below are some examples
of what could contain these variables:

\begin{verbatim}
  KANETON_USER='julien.quintard'

  KANETON_HOST='linux/ppc'
  KANETON_HOST='windows~cygwin/ia32'

  KANETON_PYTHON='/usr/bin/python'

  KANETON_PLATFORM='ibm-pc'
  KANETON_PLATFORM='sgi/o2'
  KANETON_PLATFORM='sgi/octane'

  KANETON_ARCHITECTURE='mips64'
  KANETON_ARCHITECTURE='ia32/educational'
  KANETON_ARCHITECTURE='ia32/smp'
\end{verbatim}

%
% organisation
%

\subsubsection{Organisation}

The development environment configuration files and scripts are located in
the \name{environment/} directory. The directory contains the three
following scripts:

\begin{verbatim}
  critical.py
  initialize.py
  clean.py
\end{verbatim}

The \location{critical.py} script essentially generates a configured
development environment. The result of this generation are two files called
\location{env.mk} and \location{env.py} which contains the configured
environment variables and functions for the \name{Make} files and \name{Python}
scripts, respectively. This file is called critical because it does not rely
on the portable development environment as it generates it.

The \location{initialize.py} script relies on the file \location{env.py}
previously generated. This script set up everything necessary for building the
kaneton microkernel based on the configured environment.

Finally, the \location{clean.py} script cleans everything installed by the
\location{initialize.py} script and removes the generated configured
environment files.

The generation of the configured environment is done by going through
the configuration files of all the profiles and sub-profiles associated
to the user configuration. In other words, the kaneton environment engine
processes the configuration files according to the layered organisation
described below, starting with the lowest-level layer thourgh the highest one.

\begin{verbatim}
  profile/
  profile/host
  profile/host/${KANETON_HOST}/${KANETON_ARCHITECTURE}
  profile/kaneton
  profile/kaneton/core
  profile/kaneton/platform
  profile/kaneton/platform/${KANETON_PLATFORM}
  profile/kaneton/architecture
  profile/kaneton/architecture/${KANETON_ARCHITECTURE}
  profile/kaneton/libc
  profile/user
  profile/user/${KANETON_USER}         
\end{verbatim}

In this layered organisation, a variable defined in, for instance, the
\name{host} profile could be overwritten anywhere in the upper-level layers
\location{kaneton/},
\location{kaneton/architecture/\$\{KANETON\_ARCHITECTURE\}/}, \location{user/}
\etc{}

The \name{host} and \name{kaneton} profiles are theoretically completed
separated. However, the environment engine does not check for such
unauthorised overridings. Therefore the \name{core} configuration could
override a variable previously defined in the \name{host} profile.

Finally, the \name{user} profile can override any definition to adjust the
environment to its needs.

The environment engine looks for the following types of file in the
kaneton environment profile directories:

\begin{itemize}
  \item
    \location{.conf}: the \name{configuration} files contains variable
    definitions. These files are gathered by the development environment
    engine for generating the configured environment files.
  \item
    \location{.desc}: these \name{description} files contain descriptions of
    the variables of the current profile or sub-profile. These descriptions
    are used by the \name{configure} tool.
  \item
    \location{.mk}: the \name{Make} files usually contain the implementation
    of the kaneton \name{Make} interface.
  \item
    \location{.py}: the \name{Python} files usually contain the
    implementation of the kaneton \name{Python} interface.
\end{itemize}

The engine supposes that there is no variable or function overriding in
a single profile. More precisely, if there are more than a single
configuration file in a directory, the engine cannot guarantee anything
on the order these files will be processed. As a result, the overridings
could differ depending on the processing order.

Besides, although the environment engine gathers every \name{configuration}
files it finds in the environment profiles directory, it is highly recommended
to provide a single \name{configuration} file per profile directory. This
file should be named according to the name of its profile. For more details,
take a look at the environment directory which contains existing profiles.

Moreover, the \name{configure} tool requires the \name{user} profile to
contain a single \name{configuration} file named, as explained above,
\location{\$\{KANETON\_USER\}.conf}.

The kaneton development environment engine first gathers the
\name{configuration} files and process them creating an in-memory list of
configuration variables. Then it generates the configured environment files
\location{env.mk} and \location{env.py}. Indeed, the engine outputs the
configuration variables in each file and then append the content of the
\name{Make} files and \name{Python} files of the profiles to the
configured environment files \location{env.mk} and \location{env.py},
respectively.

Note that a special rule is included in \location{env.mk} so that the
configured files are regenerated if the environment engine detects that
an environment file has been changed since the last initialization.

%
% syntaxes
%

\subsubsection{Syntaxes}

% description

\subsubsubsection{Description}

The \name{description} files describe the environment variables in order
to specify what kind of value a variable can take \etc{}

The description syntax is based on the \name{YAML} language.

Examples of variable descriptions named \code{\_FOO\_}, \code{\_BAR\_}
and \code{\_CHICHE\_} are given next:

\begin{verbatim}
  #
  # _FOO_
  #
  - variable: _FOO_
    string: the foo flag
    type: set
    values:
      Off: -D_FOO_FLAG_=0
      On: -D_FOO_FLAG_=1
    description: |
      This is a description of a two-state variable _FOO_.

  #
  # _BAR_
  #
  - variable: _BAR_
    string: the bar parameter
    type: set
    values:
      Simple: ${BAR_SIMPLE}
      Normal: ${BAR_NORMAL}
      Optimised: ${BAR_OPTIMISED}
    description: |
      This is another parameter which can take three values: simple,
      normal and optimised.

  #
  # _CHICHE_
  #
  - variable: _CHICHE_
    string: the most powerful optimisation
    type: any
    description: This is the magic kaneton optimisation.
\end{verbatim}

Note that the environment engine never takes these descriptions into account.
Indeed, this is the role of the \name{configure} tool.

In this syntax, variables are classified according to the type of value
they can take: \code{set} and \code{any}.

A \code{set} variable can take any value in a given list of values. In
this case, the \code{values} field contains couples string/value the
variable can take. The \code{string} is displayed by the \name{configure}
tool while the \code{value} is assigned to the variable.

Finally, a \code{any} variable represents a variable which can take any
value.

The \code{string} fields were introduced to avoid displaying internal
non user-friendly variable names and/or values. Therefore, the
\name{configure} tool will always display literal \code{strings}
rather than variable names or values which are likely to do not make any
sense to the user.

% configuration

\subsubsubsection{Configuration}

The \name{configuration} files contains the actual variable definitions
through a very simple syntax.

The syntax allows both assignments and completion as shown in the next example:

\begin{verbatim}
  FOO = bar
  FOO += baz
  FOO = kaneton
\end{verbatim}

The \code{FOO} variable first takes the initial value \code{bar}. Then,
the value \code{baz} is added to the previous \code{FOO}'s value
leading the the value \code{bar baz}. Finally, the last assignment
overwrites the previous definition by setting the value of \code{FOO}'s
variable to \code{kaneton}.

The configuration syntax enables variable references. These references can
be both environment variable or shell variable. The following example
illustrates this.

\begin{verbatim}
  BAR = ${FOO} is a very powerful microkernel
  SH = the shell currently used is $(SHELL)
\end{verbatim}

The reader certainly notice the \code{\$\{\}} construct is used to reference
a kaneton environment variable while the \code{\$()} one references a shell
variable.

Finally, a configuration file can also include another file using the
\code{include} statement:

\begin{verbatim}
  include ../an/other/file/far/../far/../away
\end{verbatim}

This construct is very useful to centralize definitions common to multiple
profiles or sub-profiles in a single location.

Note that kaneton environment variables start and end with an underscore
to avoid naming collisions. Another solution would have been to use a
prefix \code{KANETON\_} as it is stipulated in the general kaneton rules
but this would have led to very long names.

% make

\subsubsubsection{Make}

The \name{Make} files must implement the whole kaneton \name{Make}
interface which will be described next.

The syntax used in these files is based on the \name{GNU Make} syntax.

% python

\subsubsubsection{Python}

The \name{Python} files must implement the whole kaneton \name{Python}
interface.

The syntax used in these files is based on the \name{Python} syntax.

%
% interfaces
%

\subsubsection{Interfaces}

% make

\subsubsubsection{Make}

In this section we will detail the make interface that every host profile
must implement. The reader should look closer to the host profiles already
implemented.

Since the \name{GNU Make} syntax does not provide any name space feature,
every kaneton \name{Make} function is prefixed by \code{env\_} in order
to avoid name conflicts.

Note that the \name{Make} development environment must take care of
setting the \code{PYTHONPATH} shell environment variable with a value
including the \code{\_PYTHON\_INCLUDE\_DIR\_} kaneton environment variable
so that scripts can use kaneton \name{Python} packages.

\function{}{env\_perform}{\argument{command}}
         {
           This function performs an action according to the given
           \argument{command} argument.

           \-

           Additionally, if the \code{\_OUTPUT\_} environment variable is
           set to \code{\_OUTPUT\_VERBOSE\_}, the function displays the
           command on the output before performing it.
         }

\function{}{env\_display}{\argument{color},
                          \argument{action},
                          \argument{file},
                          \argument{indentation},
                          \argument{options}}
         {
	   This function displays a message representing an action performed
	   by the kaneton \name{Make} interface.

	   \-

	   The option \code{\$(OPTION\_NO\_NEWLINE)} can be used not to
	   output the trailing newline.
	 }

\function{}{env\_cd}{\argument{directory},
                     \argument{options}}
         {
	   This function changes the current working directory.
	 }

\function{}{env\_pull}{\argument{file},
                       \argument{options}}
         {
	   This function returns the content of the \argument{file}.
	 }

\function{}{env\_launch}{\argument{file},
                         \argument{arguments},
                         \argument{options}}
         {
	   This function launches a new program/script/make \etc{}

	   \-

	   This function must look at the file name in order to determine
	   how to launch it. Moreover, the function must move the the
	   directory where is located the file before launching it.
	 }

\function{}{env\_preprocess}{\argument{preprocessed file},
                             \argument{c file},
                             \argument{options}}
         {
	   This function launches the C preprocessor the \argument{c file}
	   and generates the \argument{preprocessed file}.
	 }

\function{}{env\_compile-c}{\argument{object file},
                            \argument{c file},
                            \argument{options}}
         {
	   This function compiles a \argument{c file} generating an
	   \argument{object file}.
	 }

\function{}{env\_lex-l}{\argument{c file},
                        \argument{lex file},
                        \argument{options}}
         {
	   This function generates a \argument{c file} from a
	   \argument{lex file}.
	 }

\function{}{env\_assemble-S}{\argument{object file},
                             \argument{S file},
                             \argument{options}}
         {
	   This function assembles an \argument{S file}.
	 }

\function{}{env\_static-library}{\argument{static library file name},
                                 \argument{object files and/or libraries},
                                 \argument{options}}
         {
	   This function builds a static library from object files.
	 }

\function{}{env\_dynamic-library}{\argument{dynamic library file name},
                                  \argument{object files and/or libraries},
                                  \argument{options}}
         {
	   This function builds a dynamic library from \argument{object
           files and/or libraries}.
	 }

\function{}{env\_executable}{\argument{executable file name},
                             \argument{object files and/or libraries},
                             \argument{layout file},
                             \argument{options}}
         {
	   This function builds a executable from object files and/or
	   libraries. The \argument{layout file} describes where to
	   place the different data: code, read-only data, stack \etc{}

	   \-

	   The option \code{\$(ENV\_OPTION\_NO\_STANDARD)} tells the function
	   not to use the operating system standard stuff: libraries, includes
	   \etc{}
	 }

\function{}{env\_archive}{\argument{archive file name},
                          \argument{object files},
                          \argument{options}}
         {
	   This function builds an archive of objects from multiple
	   \argument{object files}.
	 }

\function{}{env\_remove}{\argument{files},
                         \argument{options}}
         {
	   This function removes the files in the list.
	 }

\function{}{env\_purge}{}
         {
	   This function just cleans the current working directory from
	   unecessary files.
	 }

\function{}{env\_prototypes}{\argument{files},
                             \argument{options}}
         {
	   This function generates prototypes in relation with the given
	   \argument{files}.
	 }

\function{}{env\_headers}{\argument{files},
                          \argument{options}}
         {
	   This function generates header dependencies for the \argument{files}
	   by building a \name{Make} dependency file named
	   \code{\_DEPENDENCY\_MK\_}.

	   \-

	   The generated files \code{\_DEPENDENCY\_MK\_} are removed by
	   the environment engine when cleaning the configured development
	   environment.
	 }

\function{}{env\_version}{\argument{file}}
         {
	   This function generates a version \argument{file} from the operating
	   system's informations: user, host, date \etc{}
	 }

\function{}{env\_link}{\argument{link},
                       \argument{file},
                       \argument{options}}
         {
	   This function creates a \argument{link} to the \argument{file}.
	 }

\function{}{env\_compile-tex}{\argument{file},
                              \argument{options}}
         {
	   This function compiles the file \argument{file}.tex and
	   generates a readable document.
	 }

\function{}{env\_document}{\argument{file},
                           \argument{options}}
         {
	   This function builds a document by calling the
           \code{env\_compile-tex()} function.

	   \-

           The option \code{ENV\_OPTION\_PRIVATE} configures the document
	   by setting the {\LaTeX} definition \code{$\backslash$mode} to
	   the value \code{private}. This option was introduced to deal with
	   documents which contain information which must be kept private to
	   the students.

	   \-

	   Note that a temporary file named \code{\_DEPENDENCY\_TEX\_}
	   is created by this functions storing the {\LaTeX} definition
	   \code{$\backslash$mode}. The developer should take care of
	   removing this file in the \code{clear} \name{Make} file rule.
	 }

\function{}{env\_view}{\argument{file},
                       \argument{options}}
         {
	   This function launches a viewer for the readable document
	   generated by the function \code{env\_compile-tex()}.

	   \-

	   Note that file, as for the \code{env\_compile-tex()} function,
	   does not have any filename suffix.
	 }

% python

\subsubsubsection{Python}

In this section we will detail the kaneton \name{Python} interface that
every \name{host} profile must implement.

The \name{Python} language was designed in a portable way. For this
reason, the major part of the \name{Python} interface is implemented
by the \name{host} generic profile.

Note that the \name{Python} language provides modularity through packages.
Therefore, each \name{Python} script has to import the \code{env} package
generated by the development environment engine: \location{environment/env.py}.
Then, environment functions and variables are accessed through this package.

Below are described the functions implemented by the \code{env} package.

Note that the \name{Python} development environment must take care of
setting the \code{PYTHONPATH} shell environment variable with a value
including the \code{\_PYTHON\_INCLUDE\_DIR\_} kaneton environment variable
so that scripts can use kaneton \name{Python} packages.

\function{}{display}{\argument{header},
                     \argument{text},
                     \argument{options}}
         {
	   This function outputs some text to the screen depending on the
	   \argument{header}: \code{HEADER\_NONE}, \code{HEADER\_OK},
	   \code{HEADER\_ERROR} or \code{HEADER\_INTERACTIVE}.
	 }

\function{}{pull}{\argument{file},
                  \argument{options}}
         {
	   This function returns the content of the \argument{file}.
	 }

\function{}{push}{\argument{file},
                  \argument{content},
                  \argument{options}}
         {
	   This function writes the \argument{content} in the \argument{file}.
	 }

\function{}{temporary}{\argument{options}}
         {
	   This function creates a temporary file system object.

	   \-

	   The options \code{OPTION\_FILE} and \code{OPTION\_DIRECTORY}
	   specify which type of object to create.
	 }

\function{}{cwd}{\argument{options}}
         {
	   This function returns the path of the current working directory.
	 }

\function{}{input}{\argument{options}}
         {
	   This function waits for an input from the user.
	 }

\function{}{launch}{\argument{file},
                    \argument{arguments},
                    \argument{options}}
         {
	   This function launches a new program/script/make file \etc{}

	   \-

	   This function must look at the file name in order to determine
	   how to launch it. Moreover, the function must move the the
	   directory where is located the file before launching it.
	 }

\function{}{copy}{\argument{source},
                  \argument{destination},
                  \argument{options}}
         {
	   This function copies the file \argument{source} to
	   \argument{destination}.
	 }

\function{}{link}{\argument{source},
                  \argument{destination},
                  \argument{options}}
         {
	   This function builds a link between the file \argument{source}
	   and the file \argument{destination}.
	 }

\function{}{remove}{\argument{target},
                    \argument{options}}
         {
	   This function removes the \argument{target} which can be either
	   a file or a directory.
	 }

\function{}{list}{\argument{directory},
                  \argument{options}}
         {
	   This function lists the file system objects contains in the
	   \argument{directory}.

	   \-

	   The options \code{OPTION\_FILE} and \code{OPTION\_DIRECTORY}
	   specify which type of object to list.
	 }

\function{}{cd}{\argument{directory},
                \argument{options}}
         {
	   This function changes the current working directory to
	   \argument{directory}.
	 }

\function{}{search}{\argument{directory},
                    \argument{pattern},
                    \argument{options}}
         {
	   This function looks for files matching the given \argument{pattern}.

	   \-

	   The options \code{OPTION\_FILE} and \code{OPTION\_DIRECTORY}
	   specify which type of object to list while the
	   \code{OPTION\_RECURSIVE} option tells the function to explore
	   the whole file system sub-tree.
	 }

\function{}{pack}{\argument{directory},
                  \argument{file},
                  \argument{options}}
         {
	   This function makes an archive \argument{file} of the
	   \argument{directory}.
	 }

\function{}{unpack}{\argument{file},
                    \argument{directory},
                    \argument{options}}
         {
	   This function extracts the archive \argument{file} into the
	   \argument{directory}, if specified.
	 }

\function{}{mkdir}{\argument{directory},
                   \argument{options}}
         {
	   This function builds a new directory named \argument{directory}.
	 }

\function{}{load}{\argument{file},
                  \argument{device},
                  \argument{path},
                  \argument{options}}
         {
	   This function copies the \argument{file} on the specificed
	   \argument{device}, more precisely at the location \argument{path}.
	   The device can be virtual: an image.

	   \-

	   The options \code{OPTION\_DEVICE} and \code{OPTION\_IMAGE}
	   specify on which type of device the file must be copied.
	 }

\function{}{stamp}{\argument{options}}
         {
	   This function returns a current date.
	 }

\function{}{record}{\argument{transcript},
                    \argument{options}}
         {
	   This function starts recording a session and transcripts it
           into \argument{transcript}.
	 }

\function{}{play}{\argument{transcript},
                  \argument{options}}
         {
	   This function plays a previously recorded \argument{transcript}.
	 }

\function{}{locate}{\argument{file},
                    \argument{options}}
         {
	   This function tries to locate the program \argument{file}
	   on the system.
	 }

\function{}{path}{\argument{path},
                  \argument{options}}
         {
	   This function returns information on the given \argument{path}.

	   \-

	   The options \code{OPTION\_FILE} and \code{OPTION\_DIRECTORY}
	   specify which information the caller is interested in. The option
           \code{OPTION\_EXIST} indicates whether the \argument{path} object
           exists or not.
	 }

\function{}{info}{\argument{options}}
         {
	   This function returns information on the system.

	   \-

	   The option \code{OPTION\_CURRENT\_DIRECTORY} returns the
           sequence of characters used for accessing the current directory.
	 }


%
% ---------- header -----------------------------------------------------------
%
% project       kaneton
%
% license       kaneton
%
% file          /home/mycure/kaneton/view/book/development/control-panel.tex
%
% created       julien quintard   [sun may 20 15:22:52 2007]
% updated       julien quintard   [thu may 31 06:25:01 2007]
%

%
% ---------- control panel ----------------------------------------------------
%

\subsection{Control Panel}
\label{section:control panel}

The kaneton environment allows the developer to trigger every action from
the \textit{Make} file of the project's root directory.

These actions are listed below:

\command{init}
        {
	  This action initializes the kaneton development environment by
	  invoking the \textit{init.py} script of the \textit{environment/}
	  directory.

	  \-

	  \example{\$ make init}
	}

\command{clean}
	{
	  This action cleans the kaneton development environment.

	  \-

	  \example{\$ make clean}
	}

\command{kaneton}
	{
	  This action builds the kaneton microkernel.

	  \-

	  \example{\$ make kaneton}
	}

\command{clear}
	{
	  This action removes every compiled files.

	  \-

	  \example{\$ make clear}
	}

\command{purge}
	{
	  This action cleans directories from unecessary files.

	  \-

	  \example{\$ make purge}
	}

\command{dep}
	{
	  This action generates make files dependencies.

	  \-

	  \example{\$ make dep}
	}

\command{proto}
	{
	  This action generates C prototypes.

	  \-

	  \example{\$ make proto}
	}

\command{test}
	{
	  This action runs the test suite in order to validate the kaneton
	  whole behaviour.

	  \-

	  \example{\$ make test}
	}

\command{cheat}
	{
	  This action launches the cheat tests on students kaneton
	  implementations.

	  \-

	  \example{\$ make cheat}

	  \-

	  \example{\$ make cheat-EPITA::2006.k3}
	}

\command{build}
	{
	  This action builds the boot device.

	  \-

	  \example{\$ make build}
	}

\command{install}
	{
	  This action installs the kaneton microkernel with its dependencies:
	  configuration files, bootloader etc. on the boot device.

	  \-

	  \example{\$ make install}
	}

\command{export}
	{
	  This action builds a kaneton package.

	  \-

	  \example{\$ make export}

	  \-

	  \example{\$ make export-k3,k5}
	  \-

	  \example{\$ make export-back}
	}

\command{view}
	{
	  This action generates and displays a kaneton document.

	  \-

	  \example{\$ make view}

	  \-

	  \example{\$ make view-develop}

	  \-

	  \example{\$ make view-paper::kaneton}
	}

\command{record}
	{
	  This action records a real-time session.

	  \-

	  \example{\$ make record}

	  \-

	  \example{\$ make record-basic::test.ts}
	}

\command{play}
	{
	  This action plays a recorded session.

	  \-

	  \example{\$ make play}

	  \-

	  \example{\$ make play-session::prototypes.ts}

	  \-

	  \example{\$ make play-prototy}
	}

\command{info}
	{
	  This action displays general kaneton information.

	  \-

	  \example{\$ make info}
	}


%
% external
%

\section{External}

XXX

%
% ---------- header -----------------------------------------------------------
%
% project       kaneton
%
% license       kaneton
%
% file          /home/mycure/kaneton/view/book/development/languages.tex
%
% created       julien quintard   [fri may 25 19:52:38 2007]
% updated       julien quintard   [tue jun 12 04:25:03 2007]
%

%
% ---------- languages --------------------------------------------------------
%

\chapter{Languages}
\label{chapter:languages}

This chapter focuses on the languages in the context of the kaneton microkernel
project. Indeed, many different languages are used for the microkernel itself,
but also for the scripts and documentations around the project.

\newpage

%
% ---------- text -------------------------------------------------------------
%

% make

%
% ---------- header -----------------------------------------------------------
%
% project       kaneton
%
% license       kaneton
%
% file          /home/mycure/kaneton/view/book/development/make.tex
%
% created       julien quintard   [sat jun  2 12:34:28 2007]
% updated       julien quintard   [sun jun  3 19:22:38 2007]
%

%
% ---------- make -------------------------------------------------------------
%

\section{Make}
\label{section:make}

The coding style related to \textit{Make} files is very simple.

% naming

\subsection{Naming}

First, \textit{Make} files must be named \textbf{Makefile}, with the first
letter in upper case letter.

Rules must be named with a single word in lower case letters.

% layout

\subsection{Layout}

\textit{Make} files must be organised as sections.

\begin{itemize}
  \item
    \textbf{header}: this section contains the file header which provides
    information on the file edition: creation, last update etc.
  \item
    \textbf{dependencies}: this section contains the \textit{include}
    directives.

    \begin{verbatim}
      #
      # ---------- dependencies -----------------------------------------------
      #

      include                 ../../environment/env.mk

      [...]

      #
      # ---------- dependencies -----------------------------------------------
      #

      -include                ./$(_MAKEFILE_MK_)
    \end{verbatim}
  \item
    \textbf{directives}: this section contains \textit{Make} directives
    like \textit{.PHONY}, \textit{.SUFFIXES} etc.

    \begin{verbatim}
      #
      # ---------- directives -------------------------------------------------
      #

      .PHONY:         libs clear proto dep
    \end{verbatim}
  \item
    \textbf{variables}: this section contains the variable definitions.

    \begin{verbatim}
      #
      # ---------- variables --------------------------------------------------
      #

      SUBDIRS                 :=              arch                      \
                                              kernel                    \
                                              as                        \
                                              conf                      \
                                              id                        \
                                              region                    \
                                              segment                   \
                                              debug                     \
                                              map                       

      CORE_INCLUDE            :=              $(_CORE_INCLUDE_DIR_)/core.h

      CORE_C                  :=              core.c

      CORE_O                  :=              $(CORE_C:.c=.o)
    \end{verbatim}
  \item
    \textbf{rules}: this section contains the \textit{Make} rules.

    \begin{verbatim}
      #
      # ---------- rules ------------------------------------------------------
      #

      $(LIBSTRING_LO):        $(LIBSTRING_O)
              $(call remove,$(LIBSTRING_LO),)

              $(call archive,$(LIBSTRING_LO),$(LIBSTRING_O),)

      clear:
              $(call remove,$(LIBSTRING_O),)

              $(call remove,$(LIBSTRING_LO),)

              $(call purge,)

      proto:
              $(call prototypes,$(LIBSTRING_INCLUDE),)

      dep:
              $(call remove,$(_MAKEFILE_MK_),)

              $(call dependencies,$(LIBSTRING_C),$(_MAKEFILE_MK_),)
    \end{verbatim}
\end{itemize}

% style

\subsection{Style}

Recall that the kaneton development environment provides a \textit{Make}
interface allowing \textit{Make} files to perform operations in a portable
way. Thus, \textit{Make} files must not contain any operations specific
to any operating system.

In fact, every operation should be a call to one of the \textit{Make} functions
provided by the kaneton development environment interface.

Variables must be expressed in upper case letters while rules must be in a
single lower case letters word.

Any contributor writing a Makefile should take care of properly cleaning
the directory from temporary files. Indeed, every \textit{Make} file
provides a \textit{clear} rule which is intended to clean these temporary
files. Therefore, every developer should take care of that and verify every
file was cleaned before committing.

The kaneton project provides a tool for generating \textit{C} prototypes
through a \textit{Make} rule called \textit{prototypes}. Therefore, the
\textit{Make} file must carefully use this rule to trigger the generation
of prototypes in the right header files. For more information about the
prototypes generation, please refer to \textit{Section
\ref{section:prototypes}}.

One of the other important rules is called \textit{headers} and generated
a file containing all the dependencies related to a set of \textit{C}
source files. This generated file must be included in the \textit{Make}
so that it can use it. A \textit{Make} file inclusion is present at the
end of many \textit{Make} files for this purpose:

\begin{verbatim}
  #
  # ---------- dependencies ---------------------------------------------------
  #

  -include                ./$(_MAKEFILE_MK_)
\end{verbatim}

Finally, \textit{Make} files must resolve dependencies before performing
any other task. Therefore, a rule named \textit{dependencies} launches
every \textit{Make} file in charge of one of the dependency. Note that this
\textit{dependencies} rule is automatically launched when you try to
build a component.


% python

% XXX %
% ---------- header -----------------------------------------------------------
%
% project       kaneton
%
% license       kaneton
%
% file          /home/mycure/kaneton/view/book/development/python.tex
%
% created       julien quintard   [sat jun  2 13:06:42 2007]
% updated       julien quintard   [thu jun 14 21:52:16 2007]
%

%
% ---------- python -----------------------------------------------------------
%

\section{Python}
\label{section:python}

The \textit{Python} language is the other language for which the kaneton
development environment provides an interface for performing task in
a portable way.

Historically, kaneton scripts were written in \textit{Shell} but scripts
are pretty hard to write, especially when dealing with file modification.
Even if \textit{sed}, \textit{awk} etc. tools are very powerful, this
is not a good solution for many reasons. Therefore, the scripts were
re-written in \textit{Python} whilst a \textit{Python} development environment
interface was written.

%
% naming
%

\subsection{Naming}

File names but also function names must be expressed in lower-case letters
and a dash must be used for composite names.

Comments must also be written in lower-case letters.

Global variable names, which are often used in \textit{Python}, must be
prefixed by \textit{g\_}.

As usual, names must be expressed in a single word as long as this is possible.
Moreover, the shortcut prefixes and suffixes recommended for the \textit{C}
language also stand for \textit{Python}. For more information, take a look
at the \textit{Section \ref{section:c}}.

% environment

\subsection{Environment}

Every \textit{Python} file must rely on the development environment as it
provides a way for performing operations in a very simple and portable way.

Therefore, each \textit{Python} must start by including the configured
\textit{Python} environment module located in the \textit{environment/}
directory under the name: \textit{env.py}.

Also, since operations are provided by the development environment, note
that \textit{Python} scripts must not perform operations relying on the
\textit{sys} and/or \textit{os} modules. Others modules like \textit{re},
\textit{yaml} etc. are allowed. Moreover, the \textit{sys} module is
sometimes used - badly. If you do not know whether the usage of a module
is allowed, please ask your supervisor or directly the kaneton community
through the mailing-list for instance.

%
% layout
%

\subsection{Layout}

The \textit{Python} files must be divided into sections. Below are listed
the most common sections developers should use.

\begin{itemize}
  \item
    \textbf{header}: this section contains the well-known informations about
    the file: dates, author, license etc.
  \item
    \textbf{information}: this section describes what the \textit{Pyhton}
    do or provide if it is a package.
  \item
    \textbf{imports}: this section contains the module imports.

    \begin{verbatim}
      #
      # ---------- imports ----------------------------------------------------
      #

      import env

      import sys
      import re
    \end{verbatim}
  \item
    \textbf{globals}: this section contains the global variable declarations.

    \begin{verbatim}
      #
      # ---------- globals ----------------------------------------------------
      #

      g_directories = (`book',
                       `curriculum',
                       `exam',
                       `feedback',
                       `internship',
                       `lecture',
                       `paper')
      g_store = []
      g_document = None
      g_path = None
    \end{verbatim}
  \item
    \textbf{functions}: this section contains the actual \textit{Python}
    functions.
  \item
    \textbf{entry point}: this section contains the real entry point.

    This section was introduced in order to localise the script's entry
    point more easily. Note that this section is generally located at
    the end of the script.

    \begin{verbatim}
      #
      # ---------- entry point ------------------------------------------------
      #

      if __name__ == `__main__':
        main()
    \end{verbatim}
\end{itemize}

%
% style
%

\subsection{Style}

Every function must be preceded by a comment. This comment plays the role
of visual separator but it also describes what the function does.

This comment first contains the name of the function with two parentheses.
Then follows an empty comment line and the function description.

\begin{verbatim}
  #
  # usage()
  #
  # this function prints the documents names.
  #
  def                     usage():
    location = None
    store = None

    env.display(env.HEADER_ERROR, ``usage: view.py [document]'',
                env.OPTION_NONE)

    env.display(env.HEADER_NONE, ``'', env.OPTION_NONE)

    env.display(env.HEADER_ERROR, ``documents:'', env.OPTION_NONE)

    for store in g_store:
      for location in store:
        env.display(env.HEADER_ERROR, ``  `` + location, env.OPTION_NONE)
\end{verbatim}

The identation used in \textit{Python} script must be fixed at \textit{2}
whitespaces. This is a common kaneton rule which allows files to be
readable while respecting the limitation about the \textit{80} characters
in width.

Finally, note, that in \textit{Python}, variables do not need to be declared
before being actually used. However, kaneton people believe this lead
to unreadable code as well as to common hard-to-find bugs. For this reason,
every variable must be declared at the beginning of the function. Finally,
these variable should also be initialised.


% asm

\section*{Exercice 1 : Assembleur (3 points)}

Le code assembleur suivant est execut� sur un microprocesseur Intel
80386 en mode r�el.

Indiquez, pour chaque instruction, ce qu'elle fait.

\begin{verbatim}
MOV AX, 0x1200
MOV ES, AX
MOV BX, 0x4000
MOV AL, [ES:BX]
TEST AL, AL
JZ end
ADD AL, 42
MOV AH, 0xE
INT 0x10
\end{verbatim}


% c

%
% ---------- header -----------------------------------------------------------
%
% project       kaneton
%
% license       kaneton
%
% file          /home/mycure/kaneton/view/book/development/c.tex
%
% created       julien quintard   [fri may 25 20:24:38 2007]
% updated       julien quintard   [tue jun 12 04:25:00 2007]
%

%
% ---------- c ----------------------------------------------------------------
%

\section{C}
\label{section:c}

This section is intended to detail the kaneton \textit{C} coding style.
The specifications in this document are to be well-known by every kaneton
developer.

The kaneton coding style was introduced in order to uniformize the \textit{C}
coding styles of kaneton project contributors. Indeed, as explained in this
section, the kaneton coding style is very different from the coding styles
of other open source projects especially the \textit{GNU} style.

The key aspect of the coding style is the clarity through the coherency.
Indeed, every choice was made for the good of the project so that everyone
can very easily read the source code of the kaneton microkernel.

%
% naming
%

\subsection{Naming}

Names in kaneton must comply to the following rules.

% general

\subsubsection{General}

Entities generally belong to a higher component: manager, module, package etc.
Since the \textit{C} language does not provide any namespace feature, the
programmer must rigorously name entities by prefixing them with the name
of the higher components.

Composite names must be separated by underscores: \textit{\_}.

Moreover, entities including variables, functions, macros, types etc. must
have explicit and/or mnemonic names.

\begin{verbatim}
  #define IA32_OPTION_READ              1

  #define SET_TYPE_ARRAY                0x02

  #define MIPS64_THREAD_STACKSZ         8192
\end{verbatim}

Obviously, names can be abbreviated, but only when it allows shorter code
without a loss of meaning. The following rules must be used in this way:

\begin{itemize}
  \item
    The suffix \textbf{sz} must be used to represent a size:

    \begin{verbatim}
      #define PAGESZ                    4096

      int                               modsz;
    \end{verbatim}
  \item
    The \textbf{n} prefix must be used for variables representing a number:

    \begin{verbatim}
      int                               ntasks;
    \end{verbatim}
\end{itemize}

More generally, entities must be named with a \textit{unique} word,
excluding the namespace prefix. Short less explicit names are always prefered
on very long explicit names which are then difficult to name and use.

The following example illustrates a wrong usage of names in kaneton:

\begin{verbatim}
  t_errval              segment_dump_set(i_set              segment_set,
                                         t_uint32           number_of_segmens);
\end{verbatim}

In this example, many entities are badly named and should be replaced by
something similar to the following:

\begin{verbatim}
  t_error               segment_dump(i_set              segments,
                                     t_uint32           nsegments);
\end{verbatim}

Indeed, types are used as pre-names. Therefore, the entity name should
not overlap the type name. For example, the following example is incorrect:

\begin{verbatim}
  t_vaddr               video_addr;

  i_event               eventid;
\end{verbatim}

Developers should prefer simpler names:

\begin{verbatim}
  t_vaddr               video;

  i_event               id;
\end{verbatim}

These rules must be respected, especially in the case of functions belonging
to an interface. Indeed, since interfaces are presented to the user, they
must be easy to use and self-described. Therefore, functions as well as
arguments must be very clearly named, using a single word.

Finally, names must be expressed in English, without spelling mistakes.

% capitalization

\subsubsection{Capitalization}

Entities including variables, functions, types, structures, enumerations,
unions etc. must be expressed using lower case letters, digits and underscores
only. More precisely, entity names must be matched by the following
regular expression:

\begin{verbatim}
  [a-z][a-z0-9_]*
\end{verbatim}

This is the common rule. However, some exceptions exist.

Macro names must be entirely capitalized except for macros which create
an abstraction as it is the case for the \textit{set manager} interface. Macro
arguments, must be prefixed and suffixed by an underscore so that naming
collision is avoided.

Note that the kaneton coding style names function arguments as
\textit{arguments} and macro arguments as \textit{parameters} as these
latter ones are statically computed.

\begin{verbatim}
  #define SCHEDULER_STATE_RUN           0

  #define SEGMENT_BPT_NODESZ            4096

  #define set_reserve(_type_, _args_...)                                \
    set_reserve_##_type_(_args_)
\end{verbatim}

Comments must be written in lower case letters. Indeed, no capitalization
must be used, at all.

\begin{verbatim}
  /*
   * this function shows a segment.
   *
   * steps:
   *
   * 1) get the segment object.
   * 2) compute the type string.
   * 3) compute the perms string.
   * 4) display the entry.
   * 5) call machine dependent code.
   */

  t_error                 segment_show(i_segment                  id)
  {
    [...]
  }
\end{verbatim}

% types

\subsubsection{Types}

Type names are classified according to the group they belong to.

Firstly, structures, unions and enumerations must not be directly used.
Instead, a type must be defined. In order to indicate the entity class a type
represents, prefixes are used.

Below are listed the different prefix groups:

\begin{itemize}
  \item
    Manager types are prefixed by \textit{m\_};
  \item
    kaneton object types are prefixed by \textit{o\_};
  \item
    Dispatch interfaces are prefixed by \textit{d\_};
  \item
    Architecture-dependent types are prefixed by \textit{a\_};
  \item
    Platform-dependent types are prefixed by \textit{p\_};
  \item
    kaneton identifiers are prefixed by \textit{i\_}.
\end{itemize}

Note that group prefixes can be combined in several ways. For instance,
the prefix \textit{ao\_} represents the architecture-dependent part of
an object.

If a type does not fit into on of these groups, then it can be considered as a
general purpose type. In this case, types are prefixed by \textit{t\_}.

The example below illustrates this rule:

\begin{verbatim}
  typedef struct
  {
    i_segment                     id;

    i_as                          as;

    t_type                        type;

    t_paddr                       address;
    t_psize                       size;

    t_perms                       perms;
  }                               o_segment;
\end{verbatim}

The use of \textit{C} standard types like \textit{int}, \textit{long long}
is prohibited since there exist kaneton-specific types \textit{t\_uint32},
\textit{t\_uint64} which provide a size guarantee. Indeed, the \textit{C}
standard does not specify the size of the types \textit{int}, \textit{short}
etc. depending on the architecture.

%
% layout
%

\subsection{Layout}

The global layout of files and sections of code pertaining to the \textit{C}
preprocessor, including file inclusion and inclusion protection, must
comply to specifications detailed in the following.

Note that an \textit{Emacs} configuration file is provided in
\textit{tool/emacs}. This file contains bindkeys for generating many of
the syntaxes explained in this section. If some contributors are using
a different text editor, then it is their responsibility to follow the rules
by developing a plugin for their text editor, for instance.

% file

\subsubsection{File}

\textit{C} source and header files are composed of sections. The example
below illustrates these sections:

\begin{verbatim}
  /*
   * -------- dependencies ----------------------------------------------------
   */

  #include <core/id.h>
  #include <core/types.h>

  /*
   * -------- macros ----------------------------------------------------------
   */

  /*
   * types
   */

  #define SET_TYPE_ARRAY          0x01
  #define SET_TYPE_BPT            0x02
  #define SET_TYPE_LL             0x03
  #define SET_TYPE_PIPE           0x04
  #define SET_TYPE_STACK          0x05

  /*
   * iterator's state
   */

  #define ITERATOR_STATE_USED     0x01
  #define ITERATOR_STATE_UNUSED   0x02
\end{verbatim}

Next are listed some of the most important sections:

\begin{itemize}
  \item
    \textbf{header}: this section contains the file header which provides
    information on the file edition: creation, last update etc.

    \-

    Note that the first and last authors are also specified in this header.
    Remember that kaneton is a community driven project and therefore these
    names do not represent the main file's authors.

    \-

    This section must be present in every source and header file as it
    contains the file creation and last updates.
  \item
    \textbf{information}: this section contains a general description of
    what this file provides.

    \-

    This section must be present in every source file.

    \begin{verbatim}
      /*
       * ---------- information -----------------------------------------------
       *
       * the address space manager manages address spaces.
       *
       * an address space describes process' useable memory. each address space
       * is composed of two sets.
       *
       * the first describes the segments held by this address space, in other
       * words the physical memory.
       *
       * the latter describes the regions, the virtual areas which reference
       * some segments.
       *
       * a task can give its address space to another with as_give().
       */
    \end{verbatim}
  \item
    \textbf{dependencies}: this section contains inclusions of dependency
    files.

    \-

    This is the common section which contains the well-known \textit{\#include}
    preprocessor directives.
  \item
    \textbf{includes}: this section contains inclusion statements of additional
    files.

    \-

    This section must not be misused since there also exist a
    \textit{dependencies} section.

    \begin{verbatim}
      /*
       * ---------- dependencies ----------------------------------------------
       */

      #include <core/id.h>
      #include <core/types.h>

      /*
       * ---------- includes --------------------------------------------------
       */

      #include <core/set-array.h>
      #include <core/set-bpt.h>
      #include <core/set-ll.h>
    \end{verbatim}
  \item
    \textbf{macros}: this section contains macro definitions.

    \-

    This section must not contain macro functions as there is a section
    dedicated to these.
  \item
    \textbf{macro functions}: this section contains the macro function
    definitions which are dissociated from the basic macro definitions.

    \begin{verbatim}
      /*
       * ---------- macros ----------------------------------------------------
       */

      /*
       * options
       */

      #define SET_OPT_NONE            0x00
      #define SET_OPT_FORWARD         0x01
      #define SET_OPT_BACKWARD        0x02
      #define SET_OPT_CONTAINER       0x04
      #define SET_OPT_ALLOC           0x08
      #define SET_OPT_FREE            0x10
      #define SET_OPT_SORT            0x20
      #define SET_OPT_ORGANISE        0x40

      /*
       * ---------- macro functions -------------------------------------------
       */

      #define set_type(_type_, _id_)                                    \
        set_type_##_type_(_id_)

      #define set_reserve(_type_, _args_...)                            \
        set_reserve_##_type_(_args_)
    \end{verbatim}
  \item
    \textbf{types}: this section contains de type definitions.
  \item
    \textbf{prototypes}: this section contains the prototype definititions.

    \-

    This section is specific as the kaneton microkernel project uses a
    tool for generating prototypes from \textit{C} source file. For more
    information, please refer to \textit{Section \ref{section:prototypes}}.

    \begin{verbatim}
      /*
       * ---------- prototypes ------------------------------------------------
       *
       *      ../../core/set/set.c
       *      ../../core/set/set-array.c
       *      ../../core/set/set-ll.c
       *      ../../core/set/set-bpt.c
       *      ../../core/set/set-pipe.c
       *      ../../core/set/set-stack.c
       */

      /*
       * ../../core/set/set.c
       */

      t_error                 set_dump(void);

      t_error                 set_size(i_set                          id,
                                       t_setsz*                       size);
    \end{verbatim}
  \item
    \textbf{externs}: this section contains external declarations.

    \begin{verbatim}
      /*
       * ---------- extern ----------------------------------------------------
       */

      /*
       * the init variable, filled by the bootloader, containing in this case
       * the list of segments to mark used.
       */

      extern t_init*          init;
    \end{verbatim}
  \item
    \textbf{globals}: this section contains global variable declarations.

    \begin{verbatim}
      /*
       * ---------- globals ---------------------------------------------------
       */

      /*
       * the segment manager structure.
       */

      m_segment*              segment;
    \end{verbatim}
  \item
    \textbf{functions}: this section contains function definitions.
\end{itemize}

More specifically, \textit{C} header files are likely to contain the sections:
\textit{header}, \textit{dependencies}, \textit{macros}, \textit{types},
\textit{includes}, \textit{macro functions}, \textit{prototypes} etc. while
the \textit{C} source files should contain the sections: \textit{header},
\textit{information}, \textit{includes}, \textit{externs}, \textit{globals},
\textit{functions} etc.

The \textit{header} section is basically the file header which was described
in \textit{Chapter \ref{chapter:rules}}.

% preprocessor

\subsubsection{Preprocessor}

The preprocessor directives must appear on the first column with no
identation:

\begin{verbatim}
  #ifndef CORE_SET_H
  #define CORE_SET_H              1

  /*
   * debug
   */

  #if (DEBUG & DEBUG_SET) && defined(SET_DEBUG_TRAP)

  #define set_debug(_func_, _id_, _args_...)                              \
    fprintf(stderr, '[setd] trap: %s(%qu, %s)\n',                         \
            #_func_,                                                      \
            _id_,                                                         \
            #_args_);

  #else

  #define set_debug(_func_, _id_, _args...)

  #endif

  #endif
\end{verbatim}

All header files must be protected against multiple inclusions. The
guard macro must be named according to the location of the header file
with the suffix \textit{\_H} in order to avoid guard macro name collision.

\begin{verbatim}
  #ifndef ARCHITECTURE_IA32_SEGMENT_H
  #define ARCHITECTURE_IA32_SEGMENT_H           1

  [...]

  #endif
\end{verbatim}

Spanning macros over multiple lines is encouraged for the sake of clarity. In
such cases, escaped line breaks \textit{$\backslash$-newline} must appear
on the same column. For this purpose, tabulations must be used.

Moreover, the body of macro functions must start on a new line, idented by
\textit{2} whitespaces.

\begin{verbatim}
  /*
   * traps
   */

  #define set_trap(_func_, _id_, _args_...)                               \
    (                                                                     \
      {                                                                   \
        t_error           _r_ = ERROR_UNKNOWN;                            \
        o_set*            _set_;                                          \
                                                                          \
        set_debug(_func_, _id_, _args_);                                  \
                                                                          \
        if (set_descriptor((_id_), &_set_) == ERROR_NONE)                 \
          {                                                               \
            switch (_set_->type)                                          \
              {                                                           \
                case SET_TYPE_ARRAY:                                      \
                  _r_ = _func_##_array((_id_), ##_args_);                 \
                  break;                                                  \
                case SET_TYPE_BPT:                                        \
                  _r_ = _func_##_bpt((_id_), ##_args_);                   \
                  break;                                                  \
                case SET_TYPE_LL:                                         \
                  _r_ = _func_##_ll((_id_), ##_args_);                    \
                  break;                                                  \
                case SET_TYPE_PIPE:                                       \
                  _r_ = _func_##_pipe((_id_), ##_args_);                  \
                  break;                                                  \
                case SET_TYPE_STACK:                                      \
                  _r_ = _func_##_stack((_id_), ##_args_);                 \
                  break;                                                  \
              }                                                           \
          }                                                               \
        _r_;                                                              \
      }                                                                   \
    )
\end{verbatim}

% functions

\subsubsection{Functions}

Every function must be preceded by a comment which completely describes the
actions the function performs. Moreover, some function parts need additional
descriptions. In order not to overload the function body with heavy comments,
a kaneton-specific function organisation was introduced.

Indeed, functions body are composed of \textit{steps}, each step representing
a bunch of coherent actions. These actions are not described in the body
of the function, but rather in the comment preceding the function. Thus,
the body is clearly decomposed into steps for the sake of clarity
while expressiveness is guaranted by heavy documentation contained in the
function comment.

Therefore, function definitions must comply the organisation illustrated
below:

\begin{verbatim}
  /*
   * this function adds a set descriptor to the set container.
   *
   * steps:
   *
   * 1) if the descriptor to add is the set which will contain the set objects,
   *    the container, just put it as the set container.
   * 2) otherwise, add this object in the set container.
   */

  t_error                 set_new(o_set*                          o)
  {
    SET_ENTER(set);

    /*
     * 1)
     */

    if (o->setid == set->sets)
      {
        if ((set->container = malloc(sizeof(o_set))) == NULL)
          SET_LEAVE(set, ERROR_UNKNOWN);

        memcpy(set->container, o, sizeof(o_set));

        SET_LEAVE(set, ERROR_NONE);
      }

    /*
     * 2)
     */

    if (set_add(set->sets, o) != ERROR_NONE)
      {
        cons_msg('!', 'set: unable to add this set descriptor '
                 'to the set container\n');

        SET_LEAVE(set, ERROR_UNKNOWN);
      }

    SET_LEAVE(set, ERROR_NONE);
  }
\end{verbatim}

Note that a single blank line must be put for separating pieces of code
as it leads to a more readable source code. Moreover, a blank line must
be put before and after each step comment in the function's body.

Note that comments must comply to the template exposed above. The following
are examples of bad comments:

\begin{verbatim}
  /*
  ** bad comment
  */

  /* Bad Comment
   */

  /* bad comment */

  // bad comment
\end{verbatim}

%
% style
%

\subsection{Style}

The following sections specify various aspects of what constitutes good
programming behaviour at the language level. They cover various aspects
of the \textit{C} constructs.

% blocs

\subsubsection{Blocks}

All braces must be on their own line. This rule implies a programming style
very different from the \textit{BSD} or \textit{GNU} coding styles which
put the open brace at the end of the line.

\begin{verbatim}
  if (option & OPTION_OPTIMISED)
    {
      [...]
    }
\end{verbatim}

In addition, closing braces must appear on the same column as the
corresponding opening brace.

The text between two braces must be indented by a fixed, homogeneous amount
of whitespaces. This amount is fixed to \textit{2} whitespaces. Note that
the \textit{Emacs} default identation comply to this rule.

Moreover, the braces must also be indented by the amount of \textit{2}
spaces from the previous line. However, some exceptions exist, especially
with nested block declarations.

In \textit{C} functions, the declaration part must be separated from statements
with a single blank line. Note that when there are no declarations, there
must not be any blank line separator.

Blocks are generally composed of bunches of statements and expressions. Every
developer is welcomed to put a single blank line in order to clearly separate
bunches of pieces of code.

\begin{verbatim}
  t_error                 set_get(i_set                           setid,
                                  t_id                            id,
                                  void**                          o)
  {
    t_iterator            iterator;

    SET_ENTER(set);

    if (set_locate(setid, id, &iterator) != ERROR_NONE)
      SET_LEAVE(set, ERROR_UNKNOWN);

    if (set_object(setid, iterator, o) != ERROR_NONE)
      SET_LEAVE(set, ERROR_UNKNOWN);

    SET_LEAVE(set, ERROR_NONE);
  }
\end{verbatim}

% alignment

\subsubsection{Alignment}

Declaration identifiers must be aligned with the function name, using
tabulations only. Moreover, the declarations must be ordered according
to the length of the identifier, starting with the longer, as shown below.

\begin{verbatim}
  t_error                 segment_dump(void)
  {
    t_state               state;
    t_setsz               size;
    t_iterator            i;

    [...]
  }
\end{verbatim}

In \textit{C}, pointerness is not part of the type. However, in the kaneton
coding style, the \textit{*} pointer symbol in declarations must appear next
to the type.

Function argument lists must be broken between each argument, after the
comma. In addition, the argument identifiers must be properly aligned
together, with tabulations.

\begin{verbatim}
  t_error                 as_give(i_task                  tskid,
                                  i_as                    asid)
  {
    o_task*               from;
    o_task*               to;
    o_as*                 o;

    [...]
  }
\end{verbatim}

Structures and union fields must be aligned with the type name, using
tabulations. In addition, when declaring a structure or union type, there
must be only one field declaration per line.

\begin{verbatim}
  typedef struct
  {
    i_event                       eventid;

    t_type                        type;

    u_event_handler               handler;
  }                               o_event;
\end{verbatim}

Enumerations values must be capitalized and must appear on their own line.

When an expression, declaration, assignment etc. spans over multiple lines,
the additional lines must be indented according to the type of statement.
It is indeed the responsibility of the developer to align these lines
properly. Note that \textit{Emacs}'s alignment comply to many of the
kaneton rules.

\begin{verbatim}
  int           x = y * foo(z + pow(z, 3)) +
    a_very_very_long_function_name(y, z) + bar(baz(42, 21), y * z);
\end{verbatim}

% declarations

\subsubsection{Declarations}

There must be only one declaration per line.

External declarations must not be located in functions blocks but rather in the
global scope.

Variables may be initialised at the point of declaration. This way, bugs
detection could be improved. For this purpose, however, valid expressions
are only those composed of constants and macros.

Below are illustrated some very bad constructs:

\begin{verbatim}
  t_error                 task_current(i_task*                    tsk)
  {
    i_thread              current, next;
    o_thread*             o = thread_get(current);
    extern int            sched;

    [...]
  }
\end{verbatim}

Variables must always be declared at the beginning of a block.

% statements

\subsubsection{Statements}

A single line must not contain more than one statement. In addition, commas
must not be used on a line to separate statements.

The comma must be followed by a single space, except when they separate
arguments in functions or macro functions declarations and calls, and the
argument list spans multiple lines.

The semicolon must be followed by a newline and must not be preceded by
a whitespace, except if alone on its line.

\begin{verbatim}
  {
    int                   i;

    for (i = 0; i < 256; i++)
      ;

    [...]

    obscure_example(id, as);

    [...]
  }
\end{verbatim}

There exist exception to the above rules. For more information, please
refer to the section about \textit{Control Structures}.

Statement keywords must be followed by a single whitespace, except
those without arguments. This especially implies the \texttt{return} without
argument, like \texttt{continue} and \texttt{break}, must not be separated
from the following semicolon.

Statement keywords which take an argument must enclose the argument
between parentheses, as illustrated next.

\begin{verbatim}
  {
    i_timer              id;

    while (1)
      {
        if (id != ID_UNUSED)
          continue;

        [...]

        return (0);
      }
  }
\end{verbatim}

Finally, the use of the \texttt{goto} statement must be extremely limited.

% expressions

\subsubsection{Expressions}

All binary and ternary operators must be padded on the left and right by
one space, including assignment operators.

Prefix and suffix operators must not be padded neither on the left nor on
the right.

When necessary, padding is always done with a single whitespace.

The \verb|.| and \verb|->| operators must no be padded.

Below is an example illustrating these rules.

\begin{verbatim}
  {
    int*        p;

    x = 10 * *p + reference->value++;

    x += ((*reference).tag == 1 ? 10 : 0);
  }
\end{verbatim}

There must not be any whitespaces between the function name and the opening
parenthesis in function calls.

Expressions may span over multiple lines. When a line break occurs within
an expression, it must appear just after a binary operator, in which case the
binary operator must not be padded on the right by a whitespace.

%
% control structures
%

\subsection{Control Structures}

% general

\subsubsection{General}

Control structure keywords must be followed by a whitespace. The conditional
parts of algorithmic constructs - \texttt{if}, \textit{while}, \texttt{do},
\textit{for}, \textit{else} - must be alone on their line.

The following constructs are incorrect:

\begin{verbatim}
  {
    if(option & OPTION_SET)
      return 0;

    while (str[i]) write(1, str[i++], 1)

    if (id != ID_UNUSED) {
      foo();
    } else {
      bar();
    }

    do {
      ++x;
    } while (x < 10);
  }
\end{verbatim}

The following must be prefered:

\begin{verbatim}
  {
    if (option & OPTION_SET)
      return (0);

    while (str[i])
      write(1, str[i++], 1)

    if (id != ID_UNUSED)
      {
        foo();
      }
    else
      {
        bar();
      }

    do
      {
        ++x;
      } while (x < 10);
  }
\end{verbatim}

Finally, the conditional control structures must always use a full
comparison and avoid the use of the logical not operator: \textit{!}

For instance, the following construct is very wrong:

\begin{verbatim}
  if (!(kernel_buffer = malloc(size)))
    MESSAGE_LEAVE(message, ERROR_UNKNOWN);
\end{verbatim}

Indeed, this construct supposes the \textit{malloc()} function returns
\textit{0} when it fails. However, the \textit{malloc()} function returns
\textit{NULL} when it fails and no specification actually indicates what
is the value of the \textit{NULL} macro. Many programmers consider this
macro is always equivalent to the \textit{zero} integer value but this is
not true.

Therefore, kaneton developers must use comparison operators to avoid such
issues. In addition, using comparison operators makes the code much
more simple to read and understand:

\begin{verbatim}
  if ((kernel_buffer = malloc(size)) == NULL)
    MESSAGE_LEAVE(message, ERROR_UNKNOWN);
\end{verbatim}

% for

\subsubsection{for}

As a general exception, the \texttt{for} construct breaks many of the
previously defined rules.

Multiple statements may appear in the initial and iteration part of the
\texttt{for} structure. For this effect, commas must be used to separate
statements.

The three part of the \texttt{for} construct may span over multiple lines. The
\texttt{while} construct must be prefered to a \texttt{for} construct with
three empty parts.

The following examples are very wrong:

\begin{verbatim}
  {
    int         i;

    for (i = 0, int j = 1;
         j = j + 2, i < 10;
         i++)
      ;

    for (;;) ;
  }
\end{verbatim}

Instead, prefer:

\begin{verbatim}
  {
    int         i;
    int         j;

    for (i = 0, j = 1, j = j + 2;
         i < 10;
         j = j + 2, i++)
      ;

    for (; 1; )
      ;
  }
\end{verbatim}

Finally, single-line loops - \texttt{for} and \texttt{while} - must have
their terminating semicolon on the following line, as illustrated above.

%
% kaneton
%

\subsection{kaneton}

The following rules apply specifically to the source code of the microkernel
itself.

In kaneton, function declarations follow the same template. Indeed, the
return value must always indicate whether the function call successed or not.
The return type is therefore always the same: \textit{t\_error}.

Function arguments must be ordered, starting with input argument and
finishing with ouput arguments. Moreover, the first argument must
be the more important \textit{i.e} the kaneton object identifier or
capability on which the action is performed.

\begin{verbatim}
  t_error                 task_reserve(t_class                    class,
                                       t_behav                    behav,
                                       t_prior                    prior,
                                       i_task*                    id);
\end{verbatim}

In this example, \textit{class}, \textit{behav} and \textit{prior} are
input arguments. The function creates a new task object and returns the
identifier of this new task into the \textit{id} output argument, which
is the last of the list.

Especially note that kaneton does not have any function returning a value.
The return value is always used as an error status, thus leading to a more
uniform and coherent programming style.

Functions located in a kaneton manager must start with a call to the
\textit{ENTER()} macro function and finish with a call to the \textit{LEAVE()}
macro function. These macro functions were introduced in order to
facilitate the process of adding a statement at the beginning and/or at
the end of a function. Especially, the \texttt{return} keyword should no
longer be used.

Below is an example of such calls:

\begin{verbatim}
  /*
   * this function releases a thread from a given task.
   *
   * steps:
   *
   * 1) call the machine-dependent code.
   * 2) get the thread object.
   * 3) get the task object.
   * 4) release the thread stack if needed.
   * 5) release the thread-s object identifer.
   * 6) remove the thread from the task threads list.
   * 7) remove the thread from the threads set.
   */

  t_error                 thread_release(i_thread                 threadid)
  {
    o_task*               task;
    o_thread*             o;

    THREAD_ENTER(thread);

    /*
     * 1)
     */

    if (machine_call(thread, thread_release, threadid) != ERROR_NONE)
      THREAD_LEAVE(thread, ERROR_UNKNOWN);

    /*
     * 2)
     */

    if (thread_get(threadid, &o) != ERROR_NONE)
      THREAD_LEAVE(thread, ERROR_UNKNOWN);

    /*
     * 3)
     */

    if (task_get(o->taskid, &task) != ERROR_NONE)
      THREAD_LEAVE(thread, ERROR_UNKNOWN);

    /*
     * 4)
     */

    if (o->stack)
      {
        if (map_release(task->asid, o->stack) != ERROR_NONE)
          THREAD_LEAVE(thread, ERROR_UNKNOWN);
      }

    /*
     * 5)
     */

    if (id_release(&thread->id, o->threadid) != ERROR_NONE)
      THREAD_LEAVE(thread, ERROR_UNKNOWN);

    /*
     * 6)
     */

    if (set_remove(task->threads, threadid) != ERROR_NONE)
      THREAD_LEAVE(thread, ERROR_UNKNOWN);

    /*
     * 7)
     */

    if (set_remove(thread->threads, threadid) != ERROR_NONE)
      THREAD_LEAVE(thread, ERROR_UNKNOWN);

    THREAD_LEAVE(thread, ERROR_NONE);
  }
\end{verbatim}

In this example, the types of error are not distinguished as the function
always returns \textit{ERROR\_UNKNOWN} which is not a very good idea.

XXX pointeurs de fonctions en f\_
XXX structures en s\_


% latex

%\- au lieux de \nl
%figures in .fig
%beamer, % N) devant chaque slide,
%paper: pas de \par etc.., juste des retours a la ligne pour que le document
%latex soit aussi lisible que le document genere.

% XXX %
% ---------- header -----------------------------------------------------------
%
% project       kaneton
%
% license       kaneton
%
% file          /home/mycure/kaneton/view/book/development/latex.tex
%
% created       julien quintard   [thu jun 14 09:53:16 2007]
% updated       julien quintard   [fri may 23 20:29:54 2008]
%

%
% ---------- latex ------------------------------------------------------------
%

\section{\LaTeX}
\label{section:latex}

The kaneton documents are written in the {\LaTeX} language. These documents
are viewable through the \name{view} tool. For more information about
this tool, please refer to the \reference{Section \ref{section:view}}.

%
% naming
%

\subsection{Naming}

{\LaTeX} files must be named in lower-case letters. Moreover,
composite file names must be separated by a dash \code{-}.

%
% layout
%

\subsection{Layout}

Since {\LaTeX} is not a very readable language, every contributor
is asked to follow the rules described in order to make the documents
internal representation looks like the resulted output.

Indeed, writers are asked not to use {\LaTeX} commands specifying
layout requirements especially about indentation like \verb|\paragraph|,
\verb|\vspace| \etc{} Since the kaneton project provides templates, every
kaneton document must rely on a template which specifies the paragraph
indentation, paragraph space \etc{}

Thus, the writer willing to distinguish two paragraphs in his
document should simply put a single blank line between to two texts
in his {\LaTeX} file as shown below:

\begin{verbatim}
  This is a paragraph which is intended to explain nothing special but
  how to construct paragraphs in a very simple and readable way.

  Then, after a single blank line, this text will be considered as a new
  paragraph.
\end{verbatim}

This rule aims at making the document code as readable as possible. However,
if a vertical space is needed in some place the default paragraph indentation
does not apply, then, the \verb|\-| {\LaTeX} command must be used. This is
especially useful in \name{lecture} documents.

The example below illustrates such a \name{lecture} document and the
need of vertical indentation:

\begin{verbatim}
  \begin{frame}
    \frametitle{Description}

    The MIPS processor is a 32-bit little-endian processor.

    \-

    This processor provides \textbf{32 integer registers}, from R0 to R31.

    [...]
\end{verbatim}

Obvisouly, lines must not exceed $80$ characters in width.

The templates provided by kaneton are located in the \location{view/template/}
directory. Those include \name{book}, \name{paper}, \name{exam} \etc{}

Therefore, every document should start by including the template file.
However, since the template files need to access other files like
{\LaTeX} dependency files, additional packages \etc{} these templates need to
know where the document directory is located from the root \location{view/}
directory. Every document must therefore specify the path to the
\location{view/} directory before including the template file.

Every file should then start with something a \name{setup} section similar
to the following:

\begin{verbatim}
  %
  % ---------- setup ----------------------------------------------------------
  %

  %
  % path
  %

  \def\path{../..}

  %
  % template
  %

  %
% ---------- header -----------------------------------------------------------
%
% project       kaneton
%
% license       kaneton
%
% file          /home/mycure/kaneton/view/template/book.tex
%
% created       julien quintard   [wed may 16 18:16:50 2007]
% updated       julien quintard   [sun jun  3 12:00:37 2007]
%

%
% class
%

\documentclass[10pt,a4wide]{book}

%
% packages
%

\usepackage[english]{babel}
\usepackage[T1]{fontenc}
\usepackage{a4wide}
\usepackage{fancyheadings}
\usepackage{multicol}
\usepackage{indentfirst}
\usepackage{graphicx}
\usepackage{color}
\usepackage{xcolor}
\usepackage{verbatim}

\usepackage{aeguill}

\usepackage[Lenny]{\path/package/fncychap}

\pagestyle{fancy}

%
% -
%

\renewcommand{\-}{\vspace{\parskip}}

%
% toc
%

\newcommand{\toc}
  {
    \tableofcontents
    \setlength{\footrulewidth}{0.3pt}
    \setlength{\parindent}{0.3cm}
    \setlength{\parskip}{2ex plus 0.5ex minus 0.2ex}
  }

%
% logos
%

\newcommand{\logos}
  {
    \begin{center}
      \includegraphics[scale=0.8]{\path/logo/kaneton.pdf}
    \end{center}
  }

%
% colors
%

\definecolor{functioncolor}{rgb}{0.40,0.00,0.00}
\definecolor{commandcolor}{rgb}{0.00,0.00,0.40}
\definecolor{verbatimcolor}{rgb}{0.00,0.40,0.00}
\definecolor{noticecolor}{rgb}{0.87,0.84,0.02}

%
% function
%

\newcommand\function[3]{
  \begin{tabular}{p{0.2cm}p{13.8cm}}
  & {\color{functioncolor}\textbf{#1}}#2
  \end{tabular}

  \begin{tabular}{p{1cm}p{13cm}}
  & #3
  \end{tabular}}

%
% align
%

\newcommand\align[1]{
  \\ & \hspace{#1}}

%
% argument
%

\newcommand\argument[1]{\textit{#1}}

%
% command
%

\newcommand\command[2]{
  \begin{tabular}{p{0.2cm}p{13.8cm}}
  & {\color{commandcolor}\textbf{#1}}
  \end{tabular}

  \begin{tabular}{p{1cm}p{13cm}}
  & #2
  \end{tabular}}

%
% notice
%

\newcommand\notice[1]{
  {\color{noticecolor}\textbf{Notice}}

  \begin{tabular}{p{0.2cm}p{13.8cm}}
  & #1
  \end{tabular}}

%
% subsubsubsection
%

\newcommand\subsubsubsection[1]{\textbf{#1}}

%
% example
%

\newcommand\example[1]{
  \textit{Example:}

  \begin{tabular}{p{0.2cm}p{13.8cm}}
  & \textit{#1}
  \end{tabular}}

%
% warning XXX
%

\renewcommand{\familydefault}{\sfdefault}

%
% verbatim stuff
%

\makeatletter

\renewcommand{\verbatim@font}
  {\ttfamily\footnotesize\selectfont}

\def\verbatim@processline{
  \hskip15ex{\color{verbatimcolor}\the\verbatim@line}\par
}

\makeatother

%
% header
%

\rhead{}
\rfoot{\scriptsize{The kaneton microkernel project}}

\date{\scriptsize{\today}}

\end{verbatim}

As every other type of kaneton file, the {\LaTeX} files are composed of
sections. Below are listed some of possible sections but note that it
highly depends on the organisation of the document: multiple files, \etc{}
as well as the type of document: presentation slides, paper, book \etc{}

The best way to make things properly is to look at the existing documents.

\begin{itemize}
  \item
    \term{header}: this section contains the file header which provides
    information on the file edition: creation, last update \etc{}
  \item
    \term{setup}: this section contains the {\LaTeX} setup: path to the
    \location{view/} directory, template including, title definition \etc{}

    \begin{verbatim}
      %
      % ---------- setup ------------------------------------------------------
      %

      %
      % path
      %

      \def\path{../..}

      %
      % template
      %

      %
% ---------- header -----------------------------------------------------------
%
% project       kaneton
%
% license       kaneton
%
% file          /home/mycure/kaneton/view/template/book.tex
%
% created       julien quintard   [wed may 16 18:16:50 2007]
% updated       julien quintard   [sun jun  3 12:00:37 2007]
%

%
% class
%

\documentclass[10pt,a4wide]{book}

%
% packages
%

\usepackage[english]{babel}
\usepackage[T1]{fontenc}
\usepackage{a4wide}
\usepackage{fancyheadings}
\usepackage{multicol}
\usepackage{indentfirst}
\usepackage{graphicx}
\usepackage{color}
\usepackage{xcolor}
\usepackage{verbatim}

\usepackage{aeguill}

\usepackage[Lenny]{\path/package/fncychap}

\pagestyle{fancy}

%
% -
%

\renewcommand{\-}{\vspace{\parskip}}

%
% toc
%

\newcommand{\toc}
  {
    \tableofcontents
    \setlength{\footrulewidth}{0.3pt}
    \setlength{\parindent}{0.3cm}
    \setlength{\parskip}{2ex plus 0.5ex minus 0.2ex}
  }

%
% logos
%

\newcommand{\logos}
  {
    \begin{center}
      \includegraphics[scale=0.8]{\path/logo/kaneton.pdf}
    \end{center}
  }

%
% colors
%

\definecolor{functioncolor}{rgb}{0.40,0.00,0.00}
\definecolor{commandcolor}{rgb}{0.00,0.00,0.40}
\definecolor{verbatimcolor}{rgb}{0.00,0.40,0.00}
\definecolor{noticecolor}{rgb}{0.87,0.84,0.02}

%
% function
%

\newcommand\function[3]{
  \begin{tabular}{p{0.2cm}p{13.8cm}}
  & {\color{functioncolor}\textbf{#1}}#2
  \end{tabular}

  \begin{tabular}{p{1cm}p{13cm}}
  & #3
  \end{tabular}}

%
% align
%

\newcommand\align[1]{
  \\ & \hspace{#1}}

%
% argument
%

\newcommand\argument[1]{\textit{#1}}

%
% command
%

\newcommand\command[2]{
  \begin{tabular}{p{0.2cm}p{13.8cm}}
  & {\color{commandcolor}\textbf{#1}}
  \end{tabular}

  \begin{tabular}{p{1cm}p{13cm}}
  & #2
  \end{tabular}}

%
% notice
%

\newcommand\notice[1]{
  {\color{noticecolor}\textbf{Notice}}

  \begin{tabular}{p{0.2cm}p{13.8cm}}
  & #1
  \end{tabular}}

%
% subsubsubsection
%

\newcommand\subsubsubsection[1]{\textbf{#1}}

%
% example
%

\newcommand\example[1]{
  \textit{Example:}

  \begin{tabular}{p{0.2cm}p{13.8cm}}
  & \textit{#1}
  \end{tabular}}

%
% warning XXX
%

\renewcommand{\familydefault}{\sfdefault}

%
% verbatim stuff
%

\makeatletter

\renewcommand{\verbatim@font}
  {\ttfamily\footnotesize\selectfont}

\def\verbatim@processline{
  \hskip15ex{\color{verbatimcolor}\the\verbatim@line}\par
}

\makeatother

%
% header
%

\rhead{}
\rfoot{\scriptsize{The kaneton microkernel project}}

\date{\scriptsize{\today}}


      %
      % header
      %

      \lhead{\scriptsize{The kaneton microkernel :: development book}}

      %
      % title
      %

      \title{The kaneton microkernel :: development book
             \version
             \logo}
    \end{verbatim}
  \item
    \term{text}: this section is used when the text is relatively short
    and likely to fit in a single file. The section then contains the
    whole document's text.

    Instead, a section whose name is based on the current chapter or section
    can be used if the document is much more larger and generally splitted
    into multiple files.

    \begin{verbatim}
      %
      % ---------- latex ------------------------------------------------------
      %

      \section{\LaTeX}

      The kaneton documents are written [...]
    \end{verbatim}
\end{itemize}

%
% style
%

\subsection{Style}

First of all, comments must be written in lower-case letters. Comments
are not intended, in {\LaTeX} files, to describe what the file do but instead
play the role of visual separators.

Note that documents are either \name{public} or \name{private}. Indeed,
the {\LaTeX} definition \verb|\mode|, stored in a temporary file
included by kaneton templates, can be used to hide information which must be
kept private to the kaneton developers community like, for instance,
implementation details. For more information, take a look at the template
you are interested in or at the other equivalent kaneton documents.

Recall that the \location{view/figures/} directory contains figures related
to the kaneton microkernel documents. These figures must always be prefered
to specific figures. Moreover, writers are asked to put their figures in
this directory if the figure is general enough. Furthermore, figures must be
in the \name{FIG} format. Finally, figures must be exported into a
\name{PDF} format as it is the only format which is accepted both by
pure {\LaTeX} documents as well as \name{Beamer} documents while rendering
without quality loss.

Remember that every element must be aligned according to its parent. For
instance, in the \code{enumerate} environment, the text related to the item
must be placed on a new line, aligned by two characters from the \verb|\item|
element as shown below.

\begin{verbatim}
  \begin{itemize}[<+->]
    \item
      \textbf{R2}: return value.
    \item
      \textbf{R29}: frame pointer.
    \item
      \textbf{R30}: global variables area pointer.
  \end{itemize}
\end{verbatim}

XXX

The remaining of this section is divided into two parts: \name{Article} and
\name{Presentation} because depending on the type of document, either pure
{\LaTeX} or the \name{Beamer} extension is used.

% article

\subsubsection{Article}

A long comment separator, composed of three lines, must be used before
every new subsections, while a single line separator must be used for
subsubsections and other less important components.

\begin{verbatim}
  %
  % naming
  %

  \subsection{Naming}

  Names in kaneton must comply to the following rules.

  % general

  \subsubsection{General}
\end{verbatim}

Nevertheless, this rule also depends on the document organisation.

{\LaTeX} references must be used, via the commands \verb|\label| and
\verb|\ref|, making the reading much more easier.

% presentation

\subsubsection{Presentation}

Unlike papers, presentation documents are based on the \name{Beamer} package.
These documents are therefore composed of slides or frame in the
\name{Beamer} terms.

\name{Lecture} documents are based on such a package through the kaneton
templates, especially the \name{lecture} template.

Each frame must be preceded by a commented number indicating the slide's
number into the section of subsection.

\begin{verbatim}
  %
  % introduction
  %

  \section{Introduction}

  % 1)

  \begin{frame}
    \frametitle{Description}

    \begin{itemize}[<+->]
      \item
        About \textbf{thirty} hours course.
      \item
        Concluded by an exam.
    \end{itemize}
  \end{frame}

  % 2)

  \begin{frame}
    \frametitle{Contents}

    \begin{itemize}[<+->]
      \item
        External architecture.
      \item
        Pipeline.
      \item
        Compiler optimisations.
      \item
        Memory.
    \end{itemize}
  \end{frame}
\end{verbatim}

% commands

\subsubsection{Commands}

XXX[lister les commandes en fonction du type de document: article/presentatoin]


% shell

%deconseille fortement, preferer python

% XXX \input{shell.tex}

%XXX environment: desc, conf, text en general


%
% ---------- header -----------------------------------------------------------
%
% project       kaneton
%
% license       kaneton
%
% file          /home/mycure/kaneton/view/book/core/licenses.tex
%
% created       julien quintard   [tue may 12 18:35:31 2009]
% updated       julien quintard   [tue may 12 18:35:33 2009]
%

%
% ---------- licenses ---------------------------------------------------------
%

\chapter{Licenses}
\label{chapter:licenses}

In this chapter we will details the kaneton-related licenses.

\newpage

%
% ---------- text -------------------------------------------------------------
%

The kaneton project might be considered as an open-project since source
code is provided.

Nevertheless this is not the case as this project is used as material for
operating system courses.

Therefore, people implementing the kaneton microkernel should not
make their source code available. To avoid problems, especially students
cheating, kaneton people decided to use a kaneton-specific license
forbidding source code distribution.

The \name{kaneton license} is based on a more generic license, the
\name{pedagogical licence}.

The next sections will contain these licenses' descriptions.

%
% pedagogical license
%

\section{Pedagogical License}

\textbf{TERMS AND CONDITIONS FOR COPYING, DISTRIBUTION AND MODIFICATION}

\begin{enumerate}
  \item
    This License applies to any program or other work which contains a
    notice placed by the copyright holder saying it may be distributed
    under the terms of this License. The "Program", below, refers
    to any such program or work, and a "work based on the Program"
    means either the Program or any derivative work under copyright law:
    that is to say, a work containing the Program or a portion of it,
    either verbatim or with modifications and/or translated into another
    language. (Hereinafter, translation is included without limitation in
    the term "modification".) Each licensee is addressed as "you".
  \item
    You must not copy or distribute copies of the Program's source
    code, object code or executable form without explicit authorization
    from the maintainers.

    If you have this authorization, you must conspicuously and
    appropriately publish on each copy an appropriate copyright notice
    and disclaimer of warranty; keep intact all the notices that refer
    to this License and to the absence of any warranty; and give any
    other recipients of the Program a copy of this License along with
    the Program.
  \item
    You may modify your copy or copies of the Program or any portion of
    it, thus forming a work based on the Program, provided that you
    meet all of these conditions :

    \begin{enumerate}
      \item
	You must not publish your work without explicit authorization
	from the maintainers.
      \item
	You must send to the maintainers any work that in whole or in
	part contains or is derived from the Program or any part thereof.
      \item
	You must cause any work that you send, that in whole or in
	part contains or is derived from the Program or any part thereof,
	to be licensed as a whole under the terms of this License.
      \item
	You must cause the modified files to carry prominent notices
	stating that you changed the files and the date of any change.
      \item
	If the modified program normally reads commands interactively
	when run, you must cause it, when started running for such
	interactive use in the most ordinary way, to print or display an
	announcement including an appropriate copyright notice and a
	notice that there is no warranty (or else, saying that you
	provide a warranty) and telling the user how to view a copy
	of this License. (Exception: if the Program itself is interactive
	but does not normally print such an announcement, your work based
	on the Program is not required to print an announcement.)
    \end{enumerate}
  \item
    Access to the Program's source is granted if either:

    \begin{enumerate}
      \item
	You want to make the Program evolve
      \item
	You have pedagogical goals
    \end{enumerate}

    These requirements apply to the modified work as a whole. If
    identifiable sections of that work are not derived from the Program,
    and can be reasonably considered independent and separate works in
    themselves, then this License, and its terms, do not apply to those
    sections when you distribute them as separate works. But when you
    distribute the same sections as part of a whole which is a work based
    on the Program, the distribution of the whole must be on the terms of
    this License, whose permissions for other licensees extend to the
    entire whole, and thus to each and every part regardless of who wrote
    it.

    Thus, it is not the intent of this section to claim rights or contest
    your rights to work written entirely by you; rather, the intent is to
    exercise the right to control the distribution of derivative or
    collective works based on the Program.

    In addition, mere aggregation of another work not based on the Program
    with the Program (or with a work based on the Program) on a volume of
    a storage or distribution medium does not bring the other work under
    the scope of this License.
  \item
    You may not copy, modify, sublicense, or distribute the Program
    except as expressly provided under this License. Any attempt
    otherwise to copy, modify, sublicense or distribute the Program is
    void, and will automatically terminate your rights under this
    License. However, parties who have received copies, or rights,
    from you under this License will not have their licenses
    terminated so long as such parties remain in full compliance.
  \item
    You are not required to accept this License, since you have not
    signed it. However, nothing else grants you permission to modify
    or distribute the Program or its derivative works. These actions
    are prohibited by law if you do not accept this
    License. Therefore, by modifying or distributing the Program (or
    any work based on the Program), you indicate your acceptance of
    this License to do so, and all its terms and conditions for
    copying, distributing or modifying the Program or works based on
    it.
  \item
    Each time you redistribute the Program (or any work based on the
    Program), the recipient automatically receives a license from the
    original licensor to copy, distribute or modify the Program subject
    to these terms and conditions. You may not impose any further
    restrictions on the recipients' exercise of the rights granted
    herein. You are not responsible for enforcing compliance by third
    parties to this License.
  \item
    If the distribution and/or use of the Program is restricted in
    certain countries either by patents or by copyrighted interfaces,
    the original copyright holder who places the Program under this
    License may add an explicit geographical distribution limitation
    excluding those countries, so that distribution is permitted only
    in or among countries not thus excluded. In such case, this License
    incorporates the limitation as if written in the body of this
    License.
  \item
    We may publish revised and/or new versions of this License from
    time to time. It may evolve considering new contributors
    needs. Contact us if you have any request. Such new versions will
    be similar in spirit to the present version, but may differ in
    detail to address new problems or concerns.
  \item
    If you wish to incorporate parts of the Program into other free
    programs whose distribution conditions are different, write to the
    author to ask for permission, we sometimes make exceptions for this.
  \item
    THERE IS NO WARRANTY FOR THE PROGRAM, TO THE EXTENT PERMITTED BY
    APPLICABLE LAW. EXCEPT WHEN OTHERWISE STATED IN WRITING THE
    COPYRIGHT HOLDERS AND/OR OTHER PARTIES PROVIDE THE PROGRAM "AS IS"
    WITHOUT WARRANTY OF ANY KIND, EITHER EXPRESSED OR IMPLIED,
    INCLUDING, BUT NOT LIMITED TO, THE IMPLIED WARRANTIES OF
    MERCHANTABILITY AND FITNESS FOR A PARTICULAR PURPOSE. THE ENTIRE
    RISK AS TO THE QUALITY AND PERFORMANCE OF THE PROGRAM IS WITH
    YOU. SHOULD THE PROGRAM PROVE DEFECTIVE, YOU ASSUME THE COST OF
    ALL NECESSARY SERVICING, REPAIR OR CORRECTION.
  \item
    IN NO EVENT UNLESS REQUIRED BY APPLICABLE LAW OR AGREED TO IN
    WRITING WILL ANY COPYRIGHT HOLDER, OR ANY OTHER PARTY WHO MAY
    MODIFY AND/OR REDISTRIBUTE THE PROGRAM AS PERMITTED ABOVE, BE
    LIABLE TO YOU FOR DAMAGES, INCLUDING ANY GENERAL, SPECIAL,
    INCIDENTAL OR CONSEQUENTIAL DAMAGES ARISING OUT OF THE USE OR
    INABILITY TO USE THE PROGRAM (INCLUDING BUT NOT LIMITED TO LOSS OF
    DATA OR DATA BEING RENDERED INACCURATE OR LOSSES SUSTAINED BY YOU
    OR THIRD PARTIES OR A FAILURE OF THE PROGRAM TO OPERATE WITH ANY
    OTHER PROGRAMS), EVEN IF SUCH HOLDER OR OTHER PARTY HAS BEEN
    ADVISED OF THE POSSIBILITY OF SUCH DAMAGES.
\end{enumerate}

%
% kaneton license
%

\section{kaneton License}

\textbf{TERMS AND CONDITIONS FOR COPYING, DISTRIBUTION AND MODIFICATION}

\begin{enumerate}
  \item
    This licence is nothing more than a link to the pedagogical licence.
  \item
    Any program under the kaneton licence is in fact under the terms and
    conditions of the pedagogical licence.
\end{enumerate}

%%%
%% licence       kaneton licence
%%
%% project       kaneton
%%
%% file          /home/mycure/kaneton/view/papers/kaneton/bibliography.tex
%%
%% created       julien quintard   [mon may  8 18:35:35 2006]
%% updated       julien quintard   [mon may  8 20:38:56 2006]
%%

%
% bibliograpy
%

\chapter{Bibliography}

This chapter contains the bibliography.

%
% text
%

\begin{thebibliography}{0}
  \bibitem{AST-SCO}
    \textbf{Structured Computer Organization};
    by
    \textit{Andrew S. Tanenbaum}
  \bibitem{AST-CN}
    \textbf{Computer Networks};
    by
    \textit{Andrew S. Tanenbaum}
  \bibitem{AST-OSDI}
    \textbf{Operating Systems: Design and Implementation};
    by
    \textit{Andrew S. Tanenbaum, Albert S Woodhull}
  \bibitem{AST-MOS}
    \textbf{Modern Operating Systems};
    by
    \textit{Andrew S. Tanenbaum}
  \bibitem{AST-DOS}
    \textbf{Distributed Operating Systems};
    by
    \textit{Andrew S. Tanenbaum}
  \bibitem{AST-DSPP}
    \textbf{Distributed Systems: Principles and Paradigms};
    by
    \textit{Andrew S. Tanenbaum, Maarten van Steen}
  \bibitem{NAL-DA}
    \textbf{Distributed Algorithms};
    by
    \textit{Nancy A. Lynch}
\end{thebibliography}


\end{document}

%
% ---------- kaneton ----------------------------------------------------------
%

XXX regles generales de kaneton

$$$ ne peut on pas les mettre dans development??? $$$

aucun fichier dans le repository avec les droits executables: Makefile,
scripts etc. must be non-executable.

utiliser des mots entiers plutot que des simplifications sauf si vraiment
le mot est long.

preferez - plutot que _ pour separer sauf si le langage ne le permet pas.

si le langage ne permet pas l'utilisation de namespace ou forme du genre,
prefixez tout par le nom du module.

SURTOUT, regardez le code deja produit afin de vous inspirez de ce qui
existe deja et evitez des erreurs, c'est TRES IMPORTANT!!! si ces regles
ne sont pas suivies tout le monde perd du temps et le projet perd en
clarete. de plus a long terme, du temps devra etre perdu a developper des
outils specifiques pour verifier que vous respectez ces regles ce qui serait
vraiment completement stupide.

le nommage d'une fonctionnalite doit autant que possible se faire en un
seul et unique mot, ce mot designant au mieux la fonctionnalite. il peut bien
evidemment etre prefixe par le nom du module.

arguments en C, doivent etre simple, complet clairs car c'est l'interface
du module. o doit etre utilise opur un objet kernel correspondant au module.
tout doit etre nomme en un mot autant que possible et etre coherent par
rapport a ce qui existe deja. de plus le nom d'une variable ne doit
pas recoupe son type: t_wait waitlist c'est stupide.

XXX

%
% ---------- development ------------------------------------------------------
%

pour chaque language explique les regles: header, commentaire etc.

minuscule

-- c

-- make

-- python

-- latex

figures in .fig

-- asm

-- shell

XXX
ca ne devrait pas juste expliquer les regles mais aussi comment utiliser
les outils, syntaxes etc. genre configure/, test/, cheat/ etc.

------------------------------

introduction

source-tree

tools

  internal
  external

rules

  tools
    XXX
  languages
    c
    python
    ...

XXX latex
beamer, % N) devant chaque slide,
paper: pas de \par etc.., juste des retours a la ligne pour que le document
latex soit aussi lisible que le document genere.

XXX cheat/2007 NOTHING !!!!!!!!

pas de shebang, pas d'executable

.tar.gz pour les tarballs

-- code

eviter qu'une seule personne ne soit maitre d'un gros bout de code sans
quoi si elle part, le projet tombe a l'eau. il faut donc que tlm s'interesse
aux autres parties du code en demandant si il faut des infos autres autres
developpeurs sur cette partie.

-- tool :: ml

XXX pour la ML, balise pour retrouver facilement dans les archives.
Ne repondre qu'en relation avec le topic de la discussion. si c'est autre
chose alors creer une nouvelle discussion.

mailing-list netiquette to respect
>>>>>>> .r1166
